\documentclass[11pt,a4paper]{article}
\usepackage[T1]{fontenc}
\usepackage{lmodern}
\usepackage[margin=24mm]{geometry}
\usepackage{amsmath,amssymb,amsthm,mathtools}
\usepackage{booktabs,longtable,array,enumitem,microtype,xcolor,fancyhdr}
\usepackage{xurl}
\usepackage[colorlinks=true,linkcolor=blue!40!black,citecolor=blue!40!black,urlcolor=blue!45!black]{hyperref}
\hypersetup{pdftitle={Eliminating auxiliary witnesses in Diophantine consistency},pdfauthor={AI-assisted research prepared for Vladimir Reshetnikov},pdfsubject={Explicit polynomial equations whose unsolvability is equivalent to Con(ZFC)}}
\definecolor{ink}{HTML}{19364B}
\definecolor{muted}{HTML}{526575}
\pagestyle{fancy}\fancyhf{}
\fancyhead[L]{\small\color{muted}Diophantine consistency: new certificate reductions}
\fancyhead[R]{\small\color{muted}ProveIt research}
\fancyfoot[C]{\thepage}
\setlength{\headheight}{14pt}
\setlength{\parindent}{0pt}
\setlength{\parskip}{4pt}
\setlength{\emergencystretch}{3em}
\newtheorem{theorem}{Theorem}[section]
\newtheorem{lemma}[theorem]{Lemma}
\newtheorem{proposition}[theorem]{Proposition}
\newtheorem{corollary}[theorem]{Corollary}
\theoremstyle{remark}\newtheorem{remark}[theorem]{Remark}
\newcommand{\N}{\mathbb N}
\newcommand{\Z}{\mathbb Z}
\newcommand{\Q}{\mathbb Q}
\newcommand{\ZFC}{\mathsf{ZFC}}
\newcommand{\ZF}{\mathsf{ZF}}
\newcommand{\Con}{\operatorname{Con}}
\newcommand{\code}[1]{\nolinkurl{#1}}
\DeclareMathOperator{\ord}{ord}
\DeclareMathOperator{\lcm}{lcm}
\title{\color{ink}\textbf{Eliminating Auxiliary Witnesses\\in Diophantine Consistency}\\[7pt]
\large Polynomial comparison centers, staged index recovery,\\and a descending-binomial factorial identity}
\author{Research prepared for Vladimir Reshetnikov\\
\small AI-assisted mathematical construction and independent review}
\date{25 October 2026}
\begin{document}
\maketitle\thispagestyle{empty}
\begin{abstract}
We give three complementary improvements to explicit Diophantine equations
for the consistency of ZFC, compared with ProveIt commit
\code{1c5367452a0bd4dd7e9e4bceb1508b62dbd2ff61}, including its newly placed
packages 127--133. A direct quartic in 314 unrestricted integers has 204
quadratic residuals, 5,721 collected monomials, and coefficient height
128,446. It improves the 324-variable parent in variables, residuals, and
support, at the same height. A direct quartic in 233 nonnegative integers
has 182 residuals, 7,313 monomials, and height 86,543,666; it dominates
both the 234-variable and 233-variable natural parents in the compared
metrics. Separate polynomials in ten positive integers have complete
78-operation and 79-operation evaluators, at exact degrees 630 and 468.
Their shorter arithmetic circuits trade against larger degrees and retain
their parent's large fixed coefficients.

The main new mechanisms are a deterministic polynomial in place of a
Chinese-remainder-selected Pell comparison base; recovery of an index
from a known divisor before recovering its full denominator; and the
identity
\(F=\lfloor(F^2-1)^t/\binom{F^2-1}{t}\rfloor\iff F=t!\)
on the stated positive domain. Complete new ordinary arguments,
literal polynomial data, coefficient recipes, exact reconstruction,
and independent finite checks are supplied. The direct constructions
encode the numerical proof calculus and use neither K215 simulation nor
Jones's universal polynomial. The compact branch retains its native
calibration. The theorem is an equivalence provable in ZFC, not a proof
of consistency or global minimality, and not a new end-to-end formal
proof-assistant certificate.
\end{abstract}
\newpage
{\setlength{\parskip}{0pt}\tableofcontents}
\newpage

\section{The completed equations and the comparison boundary}
\label{sec:results}

Write \(\N=\{0,1,2,\ldots\}\). There are four completed single
polynomials, with three different optimization goals. The domain is part
of each specification.

\begin{center}\small
\begin{tabular}{@{}lrrrrr@{}}
\toprule
Equation and domain & Variables & Degree & Rows & Terms & Height\\
\midrule
Direct, all integers &314&4&204&5,721&128,446\\
Direct, nonnegative integers &233&4&182&7,313&86,543,666\\
\bottomrule
\end{tabular}
\end{center}

For the unrestricted-integer endpoint, let \(r_0,\ldots,r_{203}\) be
the literal residuals in \code{direct/equation.json}, printed in
Appendix~\ref{app:integer}. Its polynomial is
\begin{equation}\label{eq:integer-main}
 P_{\Z}(z_0,\ldots,z_{313})=
 \sum_{j=0}^{203}r_j(z_0,\ldots,z_{313})^2.
\end{equation}
For the natural-number endpoint, let \(s_0,\ldots,s_{181}\) be the
literal residuals in \code{natural/equation.json}, printed in
Appendix~\ref{app:natural}. Its polynomial is
\begin{equation}\label{eq:natural-main}
 P_{\N}(x_0,\ldots,x_{232})=
 \sum_{j=0}^{181}s_j(x_0,\ldots,x_{232})^2.
\end{equation}
Both lists are fixed finite data. Their indices do not conceal a proof
length, an unbounded list, or a variable number of quantified witnesses.
Every residual has degree at most two. The complete collected expansions
are in the adjacent \code{expanded_polynomial.json} files.

\begin{center}\small
\begin{tabular}{@{}lrrr@{}}
\toprule
Compact equation and domain & Variables & Binary operations & Exact degree\\
\midrule
\code{compact/equation78.json}, positive integers &10&78&630\\
\code{compact/equation79.json}, positive integers &10&79&468\\
\bottomrule
\end{tabular}
\end{center}
Every live binary addition, subtraction, and multiplication is counted,
including squares and multiplication by fixed numerals. Named registers
are abbreviations in a polynomial expression, not extra unknowns.
Section~\ref{sec:compact-unscaled} gives the
complete formulas and proof. The four fixed integers are supplied as
literal bytes, with an independently replayed finite reconstruction.

\begin{theorem}[Consistency calibration]\label{thm:main}
For the fixed polynomials above, ordinary mathematics formalizable in
ZFC proves
\begin{align*}
 \ZFC&\vdash
  \bigl(\neg\exists\mathbf z\in\Z^{314}\ P_{\Z}(\mathbf z)=0\bigr)
       \longleftrightarrow\Con(\ZFC),\\
 \ZFC&\vdash
  \bigl(\neg\exists\mathbf x\in\N^{233}\ P_{\N}(\mathbf x)=0\bigr)
       \longleftrightarrow\Con(\ZFC).
\end{align*}
For each of the two compact polynomials \(P\), it also proves
\[
 \bigl(\neg\exists\mathbf u\in\Z_{>0}^{10}\ P(\mathbf u)=0\bigr)
       \longleftrightarrow\Con(\ZFC).
\]
Here \(\Con(\ZFC)\) is the syntactic absence of a finite refutation in
the fixed effective ZFC calculus used by the inherited calibration.
\end{theorem}

The constructive assertion proved in each direction is that a polynomial
zero exists exactly when a finite refutation exists. No consistency
assumption is used in this equivalence. The direct branches compose
with the numerical-calculus theorem
\code{KEndpoint.c222_iff_zf_refutation} and the Choice transport
\code{ConstructibleChoice.zf_refutation_iff_zfc_refutation}.
Their statements and source identities are retained. The new arithmetic
steps are proved here; the named logical calibration is inherited, with
its proof boundary recorded explicitly in Section~\ref{sec:verification}.

\subsection{What is improved}

The repository snapshot contains more than one frontier. In particular,
the collection overview had not yet incorporated all of packages
127--133 when those packages were placed. We compared their literal
endpoint files, rather than treating the older overview as the latest
result. The most relevant direct comparisons are:

\begin{center}\small
\begin{tabular}{@{}llrrrr@{}}
\toprule
Domain & Construction & Variables & Rows & Terms & Height\\
\midrule
\(\Z\) & Package 131 &324&208&5,808&128,446\\
\(\Z\) & This work &314&204&5,721&128,446\\
\(\N\) & Package 129, recommended &234&183&9,019&87,741,134\\
\(\N\) & Package 129, denser &233&182&50,762&88,674,690\\
\(\N\) & This work &233&182&7,313&86,543,666\\
\bottomrule
\end{tabular}
\end{center}

The natural polynomial has about 18.9\% fewer monomials than the
234-variable parent and about 85.6\% fewer than its 233-variable
companion. Its height is also smaller than both. The unrestricted
polynomial removes ten variables and four residuals, with smaller
support and the same height. These are strict improvements in the
displayed direct-construction comparisons.

The compact comparison is different. Package 133 supplied
84 operations at degree 436 and 85 operations at degree 328. The new
78/630 and 79/468 endpoints save six operations each but raise degree.
Neither compact endpoint is claimed to dominate the older pair in all
metrics. Nor do the direct quartics have globally fewest variables:
the repository already contains other architectures with many fewer
variables, including quartics with very large specialized coefficients.
The classical nine-natural-variable and recent ten-integer-variable
representation results concern another measure of simplicity
\cite{jones,sun}.

\subsection{Why the three mechanisms are useful}

An auxiliary witness can cost several final variables after its
arithmetic and sign certificates are expanded. The present reductions
therefore work at the level of existence proofs before optimizing the
resulting polynomial.

\begin{enumerate}[leftmargin=*]
\item A Pell comparison base previously chosen by CRT can be an explicit
polynomial in an existing Pell trace. Its defining identity supplies
both required congruences automatically. This removes a witness choice,
not merely a different name for a coordinate.
\item To identify an index, the full error denominator need not be
available first. A large known divisor, combined with a strict index
bound, suffices. The exact index then recovers the full divisibility.
This allows the compact branch to discard its artificial scale.
\item A descending-binomial ratio places the factorial on the other
side of a floor. The resulting reciprocal CRT moduli and polynomial
interfaces have a smaller collected quartic than the earlier rising
binomial construction.
\end{enumerate}

The classical Pell identities, the numerical proof calculus, coefficient
encoding, sum-of-squares compilation, semantic-anchor idea, and existing
one-row cap idea all retain their attribution. The claims of novelty
are relative to the inspected ProveIt sources and concern the stated
new lemmas or completed combinations. A literature-wide priority theorem
is not asserted.

\section{A direct quartic in 314 unrestricted integers}
\label{sec:direct314}

This component uses the numerical C222 proof calculus. It has no K215
simulation and no specialization of a universal polynomial. Its complete
ordinary polynomial is
\begin{equation}\label{eq:direct314}
 P_{314}(z_0,\ldots,z_{313})=\sum_{j=0}^{203}r_j(z)^2=0,
 \qquad z\in\Z^{314}.
\end{equation}
Every quadratic residual is printed in the appendix and supplied literally
in \code{direct/equation.json}; every collected coefficient is in
\code{direct/expanded_polynomial.json}. The new equation has exact degree
four, 5,721 monomials and coefficient height 128,446.

\begin{theorem}\label{thm:direct314}
With the retained ordinary C222-to-ZFC refutation calibration,
\[
 \ZFC\vdash
 \bigl(\neg\exists z\in\Z^{314}\;P_{314}(z)=0\bigr)
 \ \longleftrightarrow\ \Con(\ZFC).
\]
\end{theorem}

Here ``retained calibration'' identifies previously proved mathematical
dependencies; it is not an added consistency hypothesis. The present
package does not claim a new end-to-end proof-assistant derivation.
The exact parent comparison, independently recomputed from the literal
equations, is
\begin{center}
\begin{tabular}{lrrrrr}
\toprule
Construction & Variables & Rows & Degree & Terms & Height\\
\midrule
Package 131 parent &324&208&4&5,808&128,446\\
This construction &314&204&4&5,721&128,446\\
\bottomrule
\end{tabular}
\end{center}
Thus the reduction saves ten coordinates, four quadratic rows and 87
collected monomials, without increasing degree or coefficient height.
It improves this direct-calculus comparison; other coefficient
architectures already use fewer variables.

\subsection{The arithmetic interface and the semantic anchor}

The pinned parent is package 131 at commit
\code{1c5367452a0bd4dd7e9e4bceb1508b62dbd2ff61}. Its accepted factored
quadratic source is retained as \code{direct/baseline/factored.json}, and
its complete arithmetic argument as \code{direct/proofs/parent131.md}.
The front has 18 power requests sharing 15 independent Pell pairs.
One requested exponent is the derived syntax-label polynomial
\[
 h=J(A_a,A_b,A_c,A_d).
\]
The corresponding pair is literally \code{index15.X,index15.Y}, and
its exponent is \code{1+add.380}. No new unknown is introduced by
writing $h$.

For the numerical implication constructor
\[
 I(a,b)=2(a+b)(a+b+1)+4b,
\]
$J$ is the fixed sum of the nine axiom conclusions, $a,b,I(a,b)$,
the universal quantification constructor and 1. Its original natural circuit has
nonnegative coefficients. The retained encoded-column guards give
$A_a,A_b,A_c,A_d\ge1$, so
\begin{equation}\label{eq:Janchor}
 h\ge16,\qquad h\ge I(A_a,A_b)\ge2A_a^2.
\end{equation}
These are raw circuit consequences, prior to interpreting any proof.

The parent has an additional fully indexed Pell cap at its frame index
$W$. We reuse the pair $(U,C)$ already requested at $h$ as that cap.
Its principal norm is retained; the old separate cap norm is deleted.
The index-error pool continues to include all fifteen principal pairs,
including the anchor. The new frame and cap conditions are
\begin{align}
 W&=2\left(2+\sum_{a\in\mathcal A}a^2\right),\label{eq:evenframe314}\\
 C^2-W^2-1-\sum_{i\ne *}y_i^2&\ge0.\label{eq:squarecap314}
\end{align}
Here $\mathcal A$ is precisely the retained list of all actual base,
exponent and output atoms, including $h$, with their complete affine
restorations. The asterisk marks the anchor. The second condition is
one quadratic row using
\[
 \mathcal N(a,b,c)=a^2+b^2+c(c+1).
\]
This form is nonnegative on $\Z^3$ and represents every natural integer:
apply the three-square theorem to $4n+1$, whose three square residues
force one odd and two even coordinates.

All other front signs and arithmetic requests are retained. We remove
the old sign guard on the scale $A$ and the old sign guard on the
auxiliary multiplier, with their three private witnesses each. The
ordinary sign argument is given below, rather than borrowing those
deleted hypotheses.

\subsection{An odd, large scale with either sign}

Write $X_n(b),Y_n(b)$ for the integer Pell polynomials, defined by
$X_0=1,Y_0=0$ and
\[
 X_{n+1}=bX_n+(b^2-1)Y_n,
 \qquad Y_{n+1}=X_n+bY_n.
\]
For negative bases they satisfy
\[
 X_n(-b)=(-1)^nX_n(b),\qquad
 Y_n(-b)=(-1)^{n-1}Y_n(b)\quad(n\ge1).
\]
Negative indices use $X_{-n}=X_n$ and $Y_{-n}=-Y_n$.

The original odd-ordinate frame norm remains
\begin{equation}\label{eq:framenorm314}
 A^2-W^3(W+2)(2z+1)^2=1.
\end{equation}
Equation~\eqref{eq:evenframe314} makes $W$ even and $W\ge516$, hence
$A$ is odd. Put $a=|A|$. Classification at center $W+1$ shows that
$(a,|W(2z+1)|)$ has a positive Pell index $n$. As
$Y_n(W+1)\equiv n\pmod W$, this index is a positive multiple of $W$.
Consequently
\begin{equation}\label{eq:scale314}
 a\ge X_W(W+1)>W^W,
 \qquad |A-1|\ge a-1>4W^3.
\end{equation}
Every semantic atom has absolute value below $W$. In particular
$|k_i|<W$ for every requested index and $16\le h<W$.

\subsection{A deterministic comparison center}

The crucial new local construction replaces a CRT choice by the
polynomial
\begin{equation}\label{eq:detcenter314}
 \boxed{D=A+(1-A)H^2.}
\end{equation}
The norms and auxiliary multiplier are
\begin{align}
 U^2-(A^2-1)C^2&=1,\label{eq:anchornorm314}\\
 H^2-(A^2-1)J_c^2&=1,\label{eq:commonnorm314}\\
 J_c&=C^2(2e+1),\qquad e\in\Z.\label{eq:oddaux314}
\end{align}
The cap gives $|C|>W$, so $C,J_c$ are nonzero. From the common norm,
\begin{equation}\label{eq:centeridentity314}
 D=1-(A-1)(A^2-1)J_c^2.
\end{equation}
Since $A$ is odd and $C^2\mid J_c$, equations
\eqref{eq:detcenter314}--\eqref{eq:centeridentity314} imply
\begin{equation}\label{eq:centerresidues314}
 D\equiv A\pmod H,\qquad D\equiv1\pmod{4C},\qquad |D|>1.
\end{equation}
For $A\ge3$, the second expression is less than $-1$; for $A\le-3$
it exceeds 1. Thus neither a sign condition on $D$ nor an adjustable
comparison-center parameter is needed.

Retain the comparison relations
\begin{equation}\label{eq:comparison314}
 S^2-(D^2-1)Q^2=1,
 \qquad S=U+Hg,
 \qquad Q=h+4Cq.
\end{equation}
We now establish the exact absolute anchor index.

\begin{lemma}\label{lem:anchor314}
Under the displayed conditions, $|C|=Y_h(|A|)$.
\end{lemma}
\begin{proof}
Classify the first two norms at $a=|A|$:
\[
 |U|=X_i(a),\quad |C|=Y_i(a),\quad
 |H|=X_l(a),\quad |J_c|=Y_l(a),
\]
with positive indices $i,l$. The divisibility $C^2\mid J_c$ gives
$i<l$ and $|C|\mid l$. The latter is the usual squared-ordinate lemma:
first $Y_i\mid Y_l$ implies $i\mid l$ by the addition law. Write
$l=ik$ and expand the Pell unit power modulo $Y_i^2$ to obtain
\[
 Y_l/Y_i\equiv kX_i^{k-1}\pmod{Y_i}.
\]
The norm gives $\gcd(X_i,Y_i)=1$, so $Y_i\mid k$ and hence $Y_i\mid l$.

Classify the comparison pair at $|D|$. Polynomial dependence on the
base, $D\equiv A\pmod H$, and $S\equiv U\pmod H$, with the independent
coordinate signs absorbed, imply that its index $j$ satisfies
$j\equiv\pm i\pmod{2l}$. This is the parent's abscissa-residue lemma:
reduce the unit exponent modulo $2l$, reflect about $l$, and use strict
growth to distinguish the remaining absolute traces. Signed bases
contribute only further allowed signs.

Because $D\equiv1\pmod{4C}$, one has $Y_j(D)\equiv j\pmod{4C}$.
The last comparison congruence gives $j\equiv\pm h\pmod{4C}$.
As $|C|\mid l$, these imply $i\equiv\pm h\pmod{2|C|}$. But
$1\le i\le|C|$ and $1\le h<W<|C|$, so $i=h$.
\end{proof}

The two quadratic rows
\[
 T-H^2=0,\qquad D-A-(1-A)T=0
\]
realize the deterministic center. They replace the old presentation
$D=1+4J_c(2d+1)$, $D=A+Hf$. Actual savings are measured after full
projection; a local variable subtraction alone would not establish them.

\subsection{Signed index recovery without a sign guard on the scale}

Each principal norm gives
$y_i=\epsilon_iY_{j_i}(A)$ and $x_i=\pm X_{j_i}(A)$ for some
$j_i\ge0$. For non-anchor pairs the square cap implies
$|y_i|<|C|$, hence $j_i<h$; for the anchor $j_i=h$. All are below $W$.
The unchanged pooled row is
\begin{equation}\label{eq:pool314}
 \sum_{i=1}^{15}(y_i-k_i)^2=(A-1)q_0,
 \qquad q_0\in\Z.
\end{equation}
The polynomial recurrence gives $Y_j(A)\equiv j\pmod{A-1}$ for
every signed $A$. Thus $|A-1|$ divides
\[
 E=\sum_{i=1}^{15}(\epsilon_i j_i-k_i)^2.
\]
This finite index error satisfies
$0\le E<4\cdot15W^2<4W^3<|A-1|$. Hence $E=0$, and
\begin{equation}\label{eq:canonical314}
 y_i=Y_{k_i}(A),\qquad x_i=\pm X_{k_i}(A)
\end{equation}
for all requested indices, including zero or negative ones before the
front signs have been decoded. In particular $C=Y_h(A)$ with the correct
signed-base convention. We bound $E$, not the enormous ordinate-error
sum in~\eqref{eq:pool314}.

This signed-base argument credits the repository's earlier unoriented-cap
work, retained as \code{direct/proofs/signed_scale.md}. The present
proof rebinds it to the new complete semantic source.

\subsection{Power semantics and the inherited arithmetic endpoint}

For $k>0$, $2\le n<W$ and $|m|<W$, put $T_n=2An-n^2-1$ and
\[
 L_k=X_k(A)+(n-A)Y_k(A),\qquad
 R_k=X_k(A)-(n-A)Y_k(A).
\]
The identities
\[
 L_kR_k=1+T_nY_k(A)^2,\qquad
 L_{k+1}=nL_k+T_nY_k(A)
\]
hold for either sign of $A$, so $L_k\equiv n^k\pmod{T_n}$. The
frame gives $|T_n|>W^W$, larger than both $|m-n^k|$ and $|mn^k+1|$.
The genuine request therefore gives $m=n^k$ on the canonical abscissa
branch. The opposite branch would give $mn^k=-1$, impossible in integers.
Canonical orientation here need not mean a positive coordinate.

The initial retained rows give
$B=r+s+1\ge1$, $q=B+o+1\ge2$, $x\ge0$ and
$t=16+h_0^2\ge16$. The $q^t$ request first selects the canonical
$t$-pair. Its second request gives $V_0=(F^2+1)^t$, including the
possible base-one case. The ceiling $FB=V_0+s>0$ yields $F\ge1$,
then $B>F$ and $q\ge3$. The cofinal equation
$y=V_0(V_0-1)$ yields $y\ge F+1$. Retained encoded-column guards
then make every requested index positive, including
$K=2(2+t+x+y+F)$. The $B^K$ and $(BH_0)^K$ requests recover
\[
 N=B^K,
 \qquad q=(BH_0)^K,
 \qquad H_0=1+\sum \text{the nine positive representatives}.
\]
All remaining pairs support requests at $q$ or $q+1$, and the polynomial
addition law supplies the derived $t+1$ pair. Thus all 18 power
interfaces are genuine. Each retained interval-product guard, together
with $r_i+s_i=L_i-1\ge0$, then restores both remainder signs.

At this point the parent's ordinary arithmetic theorem applies to the
unchanged front. Its exact proof is retained, with the following chain
made explicit. The ceiling gives $V_0\le BF<B^2\le q$, so the factorial
digit is exact: $B=\binom{F^2+t}{t}$. The parent's fixed-point theorem
at $t\ge16$ yields $F=t!$. The same $B$ equals the product of the
pairwise coprime integers $d_i=1+F^2/i$, $1\le i\le t$. Amplified
prime-power ranges at $N=B^K$, with the self-sized radix $NH_0^K$,
recover every bounded local rule and derived-label constraint. The
Horn-flow multiplicity argument then gives a C222 refutation, even for
a supplied multiset containing cycles. The named logical interfaces
\code{KEndpoint.c222_iff_zf_refutation} and
\code{ConstructibleChoice.zf_refutation_iff_zfc_refutation}
give a finite ZFC refutation. The older full formal development is
identified rather than recompiled in this release.

\subsection{Complete witness construction}

Conversely, begin with a finite C222 refutation. Use the parent's
cofinal horizon $t=16+h_0^2$, truth padding, factorial, beta code and
front construction. Fix $R_0,K,B,N$ and choose the nine positive CRT
representatives modulo $d_i^{2K}$. Their common period is $N^2$.
Increase only $A_a$ by a sufficiently large multiple of $N^2$ so that
\begin{equation}\label{eq:dominate314}
 h\ge2\max_{i\ne *}k_i+1.
\end{equation}
This is possible by~\eqref{eq:Janchor}: the other varying target is
only $A_a$ itself, whereas $h\ge2A_a^2$. Every local polynomial
congruence is preserved at the original stronger precision. The parent
has no aggregate upper bound on the representatives.

Only after this lift compute $H_0$, $q=(BH_0)^K$, all power outputs,
digits and finally the doubled frame $W$. The uniform estimate
$J\le H_0^{2^{22}}$ and $2^{22}(t+1)<K$ preserve all digit bounds;
they do not require the representatives to lie below $N^2$.
This ordering avoids circular witness choices.

Choose $A=X_W(W+1)>0$, which is odd because $W$ is even.
The recurrence gives $Y_W(W+1)\equiv W\pmod{2W}$, so
\[
 z=\frac{Y_W(W+1)/W-1}{2}\in\Z
\]
satisfies the odd frame. Choose all principal pairs canonically, with
$U=X_h(A),C=Y_h(A)$. If $k_*$ is the largest other requested index,
Pell doubling and growth give
\[
 Y_{2k_*+1}(A)>(4A-2)Y_{k_*}(A)^2.
\]
Together with~\eqref{eq:dominate314} and $A>W^W$, this implies
$C>W+\sum_{i\ne *}Y_{k_i}(A)^2$. The square-cap slack is therefore
nonnegative and has a ternary representation. The pooled quotient
exists because $Y_{k_i}(A)\equiv k_i\pmod{A-1}$.

For the auxiliary pair choose $l=hC$ and
\[
 H=X_l(A),\qquad J_c=Y_l(A)=C Y_C(U).
\]
The unit-power expansion gives an integer $T$ such that
\[
 Y_C(U)=CU^{C-1}+(U^2-1)T,
 \qquad
 \frac{J_c}{C^2}=U^{C-1}+(A^2-1)CT.
\]
The quotient is a positive odd integer, since $A,U$ are odd. Thus
$e=(J_c/C^2-1)/2$ supplies~\eqref{eq:oddaux314} for every semantic
index $h$. The even frame extends package 132's even-index odd-quotient
construction to this arbitrary anchor index.

Define $D$ by~\eqref{eq:detcenter314} and take
$S=X_h(D),Q=Y_h(D)$. These signed integers satisfy the comparison
norm, and polynomial base transfer through~\eqref{eq:centerresidues314}
gives both required congruences. Their unrestricted quotient witnesses
therefore exist. Fill the unchanged arithmetic gates, front guards and
power quotients. This constructs a zero of the complete new source.

The two existential directions prove the new arithmetic theorem. They
may replace old private Pell witnesses; no pointwise bijection fixing
those witnesses is asserted.

\subsection{Exact compilation and checked scope}

The offline rebuild starts at the pinned, literal, already factored parent
source. No new Horn factorization is performed. The source certificate
lists all twelve modified or removed rows and both appended center rows.
An independently written checker reconstructs the whole changed source,
including every frame atom and all fifteen pool summands.

Integral monic projection then restores all 216 eliminated coordinates
by exact integer polynomials, checks all 94 genuinely absent free
coordinates, and verifies the complete retained maps. The final 62
coordinate translations are integral bijections, substituted into every
row. Every coefficient of the sum of 204 squares is independently
collected. Both the final and parent metrics are computed from literal
equations. Three altered-artifact controls are rejected.

These checks import no producer, optimizer or translation helper. The
separate rebuild regenerates nine mathematical files byte for byte in a
temporary directory. They establish finite algebraic claims; the
unbounded Pell, representation, arithmetic and syntactic claims are
the ordinary proofs above and in the retained dependencies.

\section{A descending-binomial quartic in 233 natural variables}
\label{sec:natural}

This branch starts from the direct numerical-calculus construction of
package 129 \cite{parent129}. It changes the factorial certificate and
the associated reciprocal CRT moduli, combines the semantic anchor with
a single cap equation, and uses a checked horizon floor of 12. It
retains the Cantor-coded calculus, six exact power requests, eleven
modular binomial requests, and the dominant denominator. Its final
polynomial has 233 variables, 182 residuals, 7,313 collected monomials,
and height 86,543,666.

\subsection{A short factorial fixed point}

The rising-binomial parent used a ceiling and had to exclude spurious
small values before obtaining a strong lower bound for the factorial
witness. Reversing the product gives a particularly direct lower bound.

\begin{lemma}[Descending factorial identity]\label{lem:floor}
Let \(t\ge3\), \(F\ge2\), and \(F^2-1\ge t\) be integers. Then
\[
 F=\left\lfloor\frac{(F^2-1)^t}{\binom{F^2-1}{t}}\right\rfloor
 \quad\Longleftrightarrow\quad F=t!.
\]
\end{lemma}

\begin{proof}
Put \(n=F^2-1\), \(T=t!\), and \(S=t(t-1)/2\). The ratio is
\[
 H=T\prod_{j=0}^{t-1}\left(1-\frac{j}{n}\right)^{-1}\ge T.
\]
If its floor is \(F\), then \(F\ge T\). For \(t\ge3\),
\(T-1>S\); this is immediate at \(t=3\) and persists by induction.
Consequently
\[
 n\ge T^2-1=(T+1)(T-1)>(T+1)S.
\]
The elementary product inequality
\(\prod_j(1-a_j)\ge1-\sum_j a_j\), for \(0\le a_j\le1\), gives
\[
 T\le H\le\frac{Tn}{n-S}
       =T+\frac{TS}{n-S}<T+1.
\]
Thus the floor is \(T\), so \(F=T\). Conversely, if \(F=T\),
the same displayed inequalities apply and the floor is \(F\).
\end{proof}

The actual polynomial graph makes \(F=f_0+2\), with \(f_0\in\N\),
and evaluates
\[
 n=(f_0+1)(f_0+3)=F^2-1.
\]
This expression has nonnegative coefficients. Its two floor rows are
\begin{equation}\label{eq:floor-rows}
 V_0=BF+r,\qquad B=r+s+1,\qquad r,s\in\N.
\end{equation}
After the arithmetic interfaces have recovered
\(V_0=n^t\) and \(B=\binom nt\), these rows say exactly
\(BF\le V_0<B(F+1)\). The literal shift \(B=b_0+2\) ensures
\(B\ge2\), so the recovered binomial is nonzero and \(n\ge t\).
Lemma~\ref{lem:floor} then gives \(F=t!\). No factorial, floor, or
binomial primitive survives in the final polynomial.

\subsection{Descending reciprocal moduli}

\begin{lemma}\label{lem:descending-crt}
Let \(F=t!\) with \(t\ge3\). For \(1\le i\le t\), put
\[
 d_i=\frac{F^2}{i}-1.
\]
These are positive integers greater than one, are pairwise coprime,
have no prime divisor at most \(t\), and satisfy
\[
 \prod_{i=1}^{t}d_i=\binom{F^2-1}{t}=B.
\]
\end{lemma}

\begin{proof}
Each \(i\mid F\), hence \(F^2/i\) is a multiple of \(F\) and
\(d_i\equiv-1\pmod F\). Every prime at most \(t\) divides \(F\),
so no such prime divides \(d_i\). If a prime divided both \(d_i\)
and \(d_j\), it would divide
\(i d_i-jd_j=j-i\), which is impossible for distinct indices in
\(\{1,\ldots,t\}\). The product identity is
\[
 \prod_{i=1}^{t}\frac{F^2-i}{i}
 =\frac{(F^2-1)\cdots(F^2-t)}{t!}.
\]
\end{proof}

The encoded local-row index \(C\) now uses the polynomial equality
\begin{equation}\label{eq:descending-index}
 C+1=F^2+Bq_C,\qquad q_C\in\N,
\end{equation}
in place of the parent's \(C+F^2+1=Bq_C\). Modulo \(d_i\),
\(F^2\equiv i\), so the recovered local index is still \(i-1\).
The quotient's natural domain is legitimate in completeness. Indeed,
\(B>F^2\): with \(n=F^2-1\), we have
\(\binom nt\ge\binom n2>n+1\) in the present range. Every positive
CRT representative of \(F^2-1\pmod B\) is therefore at least
\(F^2-1\), and \(q_C\ge0\).

An additional consequence is \(d_i\equiv-1\pmod4\), because
\(4\mid F\) in the actual \(t\ge12\) construction. Thus each
\(d_i\) has a prime divisor congruent to 3 modulo 4 with odd
multiplicity. This observation is not used to claim another compression
in the delivered equation; the preceding elementary CRT facts suffice.

\subsection{One cap equation at a semantic index}

For \(A\ge2\), write
\[
 X_j(A)+Y_j(A)\sqrt{A^2-1}=(A+\sqrt{A^2-1})^j.
\]
The nonnegative solutions of
\(x^2-(A^2-1)y^2=1\) are exactly these pairs, with unique \(j\ge0\).
The elementary recurrences give
\[
 Y_j(A)\equiv j\pmod{A-1},\quad Y_j(A)\ge j,\quad
 Y_{j+1}(A)>(2A-1)Y_j(A)\quad(j\ge1).
\]
The complete one-index certificate of the parent is retained for a
single pair \((X_*,Y_*)=(X_H(A),Y_H(A))\). The new simultaneous
constraint is
\begin{equation}\label{eq:natural-one-row}
 \boxed{\quad
 Y_*=H+\sum_{i=1}^{r}(y_i-k_i)^2+(A-1)Q,
 \qquad Q\in\N.\quad}
\end{equation}
It combines the size cap and the error congruence in one quadratic row.
Before the anchor is interpreted, this row already gives \(H\le Y_*\).
It therefore supplies the index bound required by the retained
one-index certificate.
This is a specialization and combination of the earlier one-row cap
idea \cite{cap128} with package 129's semantic anchor, not a priority
claim for the general idea of combining a cap and a congruence.

\begin{lemma}\label{lem:natural-cap}
Assume \(W\ge3\), \(2\le H<W\), \(r\le W-2\),
\(0\le k_i<W\), \(A>W+1\), and \(A-1>W^3\).
Let \(Y_*=Y_H(A)\), and let each \((x_i,y_i)\) be a
nonnegative Pell solution at base \(A\). If
\eqref{eq:natural-one-row} holds, then
\((x_i,y_i)=(X_{k_i}(A),Y_{k_i}(A))\) for every \(i\).
Conversely, for the positive indices in this construction, the
canonical pairs satisfy \eqref{eq:natural-one-row} with a natural
\(Q\) whenever \(H>2\max_i k_i\).
\end{lemma}

\begin{proof}
Write \(y_i=Y_{j_i}(A)\). If \(j_i\ge H\), then
\(y_i\ge Y_*\ge Y_2(A)=2A>2W\), so
\[
 (y_i-k_i)^2\ge(Y_*-W+1)^2>Y_*.
\]
This contradicts the nonnegative right side of
\eqref{eq:natural-one-row}. Therefore \(j_i<H<W\).
Reducing the row modulo \(A-1\) gives divisibility of
\[
 E=\sum_i(j_i-k_i)^2,
 \qquad 0\le E\le r(W-1)^2<W^3<A-1.
\]
Hence \(E=0\), proving every index equality.

For completeness put \(m=\max_i k_i\ge1\). The identity
\[
 Y_{2m+1}(A)=Y_{m+1}(A)^2-Y_m(A)^2
            >4A(A-1)Y_m(A)^2
\]
and \(H\ge2m+1\) show
\[
 Y_H(A)-H-\sum_i(Y_{k_i}(A)-k_i)^2
 >[4A(A-1)-r]Y_m(A)^2-H>0.
\]
This integer is divisible by \(A-1\) by the recurrence congruence.
Its quotient is the required nonnegative \(Q\).
\end{proof}

The anchor is the existing derived label polynomial \(H=J\), which
includes the Cantor constructor
\[
 I(a,b)=2(a+b)(a+b+1)+4b,
 \qquad J(a,b,c,d)\ge I(a,b)\ge2a^2.
\]
As in the parent, the nine independent CRT representatives initially
lie in \([1,N^2]\). Increase the encoded \(a\)-column by \(N^2\).
All local polynomial congruences at precision \(d_i^{2K}\) are
preserved, and their aggregate remains at most \(10N^2\).
The new \(a>N^2\) gives \(H>2N^4\). Every other required index
is at most \(2N^2\), so \(H>2\max_i k_i\), as required for the
new converse. No extra dominance inequality appears in the polynomial.
Soundness uses only the cap lemma's explicit frame hypotheses.

\subsection{The six exact powers and the order of recovery}

The semantic parameters are
\[
 t=t_0+12,\quad F=f_0+2,\quad
 R=2+t+x+y+F,\quad K=2R,\quad B=b_0+2.
\]
The six exact requests are
\begin{equation}\label{eq:natural-six}
 V_0=(F^2-1)^t,\qquad N=B^K,\qquad q=N^K,\qquad
 D_t=q^t,\quad D_y=q^{y+1},\quad D_x=q^{x+1}.
\end{equation}
These are meanings proved for finite polynomial graphs, not operators
allowed in \(P_{\N}\). The extra denominator \(q^{t+1}=qD_t\)
uses a multiplication already accounted for in the graph.

The unchanged enlarged growth frame is
\[
 W=16+\sum_{a\in\mathcal A}a,\quad Z=z+1,\quad
 L^2-\big((W+1)^2-1\big)(WZ)^2=1,\quad A=L+\sigma,
\]
where \(\mathcal A\) contains every distinct base, exponent, and
proposed output in \eqref{eq:natural-six}, as well as every modular
upper index. All these atoms are natural and less than \(W\).
Pell classification at base \(W+1\), followed by reduction modulo
\(W\), shows that its nonzero index is at least \(W\). Therefore
\[
 A\ge L\ge X_W(W+1)\ge(W+1)^W,
 \qquad A-1>W^3.
\]
The actual base \(A\) remains free to increase after the lower witness
\(L\) is chosen. This freedom is needed for the later root congruence.

There are fourteen ordinary index pairs and one fully certified
anchor. For ordinary pairs the graph uses consecutive ordinates
\(y,z\) and a product wire \(u\):
\[
 u=yz,\qquad y^2+z^2-2Au=1.
\]
Every such pair has an output row of the form
\(ny=z+R_0\), with \(R_0\ge0\) and \(1\le n<A\).
Thus \(y>0\), \(z<Ay\), and \(x=Ay-z>0\) gives an ordinary
Pell pair. The one-row cap authenticates all its indices before
interpreting a power or binomial coefficient.

For a proposed power \(m=n^k\), put
\(T_n=2An-n^2-1\) and impose
\[
 X_k(A)+(n-A)Y_k(A)=m+T_n h_n,\qquad h_n\in\N.
\]
Modulo \(T_n\), the left side is \(n^k\). The frame proves
\(T_n>W^W+W\), \(m<W\), and \(n^k<W^W\), so both small
residues must agree. This proves all six exact powers. The canonical
left side is positive and \(0\le m<T_n\), so their difference cannot
be a negative multiple of \(T_n\). Its congruence quotient is therefore
natural in completeness. The literal modulus row also excludes a zero
power base.

Before using any factorial semantics, \(B\ge2\) gives
\[
 q=B^{K^2}\ge2^{4R^2}.
\]
With \(n=F^2-1<R^2\), every coefficient \(\binom nj\),
\(0\le j\le t\), is less than
\(R^{2t}\le2^{2Rt}<2^{2R^2}<q\). Thus the first binomial digit
can be decoded before \(F=t!\) is known. Its output is \(B\);
the floor rows now recover \(F=t!\) by Lemma~\ref{lem:floor}.
There is no circular assumption that a large proposed factorial is
already the true factorial.

\subsection{The common root and exact digits}

The unchanged endpoint rows include
\[
 y=F+1+r_0,\quad
 x=F+1+(y+1)u_0,\quad
 x=1+\bigl(1+(t+1)y\bigr)u_t.
\]
They imply \(x\ge2+(t+1)y>\max(t+1,y)\). Consequently
\(M=ND_x=Nq^{x+1}\) is a multiple of every needed digit modulus.
Put \(b=q+1\) and choose the common-root equation
\[
 2Ab=b^2+1+M h_0.
\]
It implies
\[
 \rho_n=X_n(A)+(b-A)Y_n(A)\equiv b^n\pmod M.
\]
For an exact binomial output \(C=\binom nk\), with \(d=q^k\),
the literal rows are
\[
 \rho_n=r+dC+qd h,\qquad d=r+s+1,\qquad q=C+o+1.
\]
For a range request \(N\mid\binom nk\), they are
\[
 \rho_n=r+Nd h,\qquad d=r+s+1.
\]
The carry-free binomial theorem gives precisely these interpretations.
For completeness, \(\rho_n>0\), while \(0\le r+dC<qd\) for an
exact digit and \(0\le r<d\le Nd\) for a range digit. Subtracting
the indicated small residue therefore cannot give a negative multiple
of the corresponding modulus, so the reconstructed high-digit
quotients are nonnegative. Also \(N\mid q\), so \(\gcd(b,M)=1\).
The linear congruence \(2bA\equiv b^2+1\pmod M\) is soluble:
its only possible common divisor is 2, and then \(b^2+1\) is even.
There are arbitrarily large suitable bases \(A\ge L\).

The remaining ten range binomials recover the eight local source
columns, the row index, and the derived label sum. Their upper-index
and carry bounds are the parent's bounds with the descending factorial
input, which is smaller than the former \(F^2+t\). In completeness,
the semantic-anchor lift still uses the same \(10N^2\) aggregate.
For representatives \(n\le10N^2\), relevant lower indices
\(k\le R-1\), and \(0\le j\le k\), the old estimate gives
\(\binom nj<q\).
For the derived label sum, the exact source certificate gives
\(J\le(n+1)^{\ell}\), \(\ell=2^{22}\), and
\(3\ell(t+1)\le2F\le K\) gives the required stronger bound.

\subsection{Why horizon 12 is sufficient}

The literal majorant certificate has local exponent \(E=2^{23}\)
and label exponent \(\ell=2^{22}\). At the new floor,
\[
 12!=479001600,\qquad a=\lfloor\log_2(12!)\rfloor=28,
\]
and direct integer arithmetic verifies
\[
 aE<2^a,\qquad 2E\le2^a,\qquad
 3\ell\cdot13\le2\cdot12!.
\]
These imply the same unbounded inequalities used by the source.
In particular, \(R^E<2^R\) starts at \(R=2^a\); the ratio of
successive left sides is below 2 for \(R\ge2E\), so it persists
for larger integers. The label inequality persists because a factorial
grows faster than its linear comparison. Thus the smaller literal
horizon changes coefficients without weakening a soundness estimate.

\subsection{From the local arithmetic rows to refutations}

For clarity, the inherited CRT and logical steps are stated explicitly.
If \(p^e\mid\binom ak\), then \(a\bmod p^e<k\). For
\(a\ge k\), this follows from the valuation formula for binomials:
if the residue were at least \(k\), the valuation contributions at
levels \(e,e+1,\ldots\) would vanish, leaving valuation at most
\(e-1\). The case \(a<k\) is immediate. The converse is not
assumed.

For each \(d_i\), choose a prime divisor \(p\), let
\(m=v_p(d_i)\), and work modulo \(p^{Km}\). The index congruence
\eqref{eq:descending-index} and its range binomial recover \(i-1\).
The other range binomials recover the bounded local numerical values.
The source majorant bounds its nonnegative local rule polynomial by
\(R^E<2^R<2^K\). Its divisibility by \(p^{Km}\ge2^K\)
therefore makes it zero. This produces every local row of the
numerical proof calculus.

In completeness use CRT at precision \(d_i^{2K}\), whose common
period is \(B^{2K}=N^2\). A falling-factorial numerator contains a
factor divisible by \(p^{2Km}\), while division by \(k!\) loses
at most \(k-1\le K-1\) powers of \(p\). At least \(Km\)
powers remain, so all range divisibilities hold. This argument needs
only the coprimality and product of the \(d_i\), both proved for the
descending moduli above. The semantic lift preserves the same precision.

Finally the local Horn-flow construction uses pairwise coprime weights
\(w(j)=1+jF\). Its beta-coded accumulator starts at \(w(1)=F+1\)
and ends at 1. Multiplying the local flow equations and taking prime
valuations recovers the multiplicity balance of conclusions and
premises. If contradiction label 1 were outside the closure of the
selected rules, every rule concluding outside the closure would have
an outside premise; the required excess of one outside conclusion is
impossible. Conversely a finite proof tree supplies the flow, with
padding to any sufficiently large horizon. This is the retained
numerical-calculus equivalence, followed by the two logical interfaces
stated in Theorem~\ref{thm:main}.

The order is important: frame and Pell indices; six exact powers;
first exact binomial; factorial; remaining binomial ranges; finite
local rows; refutation. Completeness chooses a finite proof first,
then its bounded semantic data, then the CRT representatives and their
lift, and only afterwards the much larger Pell witnesses.

\subsection{Literal compilation and independent review}

The finite graph exports to 339 natural coordinates and 288 quadratic
rows after affine normalization. A further 106 certified eliminations
leave the displayed 233 coordinates and 182 rows. Each inverse
substitution is an integer polynomial proved nonnegative on the natural
orthant; the two restorations with negative coefficients have exact
rational sum-of-squares certificates. Expansion gives exactly 7,313
monomials and degree four. All 233 variables occur.

An independent arithmetic review checks the unbounded floor, CRT, cap,
and horizon arguments. Its separate coefficient replay verifies the
entire expanded dictionary. Bounded controls exercise 2,197 factorial
cases, ten reciprocal products, and 2,452 cap tuples; they supplement
the proofs and are not used as evidence of universal unsolvability.
The files \code{natural/REVIEW.md}, \code{natural/INDEPENDENT_REPLAY.json},
and \code{natural/REPLAY.json} give the exact scopes.


\section{An existing Pell denominator removes the artificial scale}
\label{sec:compact-unscaled}

The compact branch gives two further fixed equations in ten strictly
positive integers. Their exact arithmetic cost and total degree are
\[
 (78,630)\qquad\text{and}\qquad(79,468).
\]
Here an operation is one live binary addition, subtraction, or
multiplication, including a square or multiplication by a fixed integer.
The complete arithmetic circuits are \code{compact/equation78.json} and
\code{compact/equation79.json}. These are cost--degree tradeoffs against
the pinned package 133's $(84,436)$ and $(85,328)$ equations. The number
of unknowns and the entire fixed coefficient tuple are unchanged.

The new idea is an order of deduction. A large known divisor of a
denominator first identifies a Pell index. The identified index then
restores divisibility by the whole denominator. Pairwise coprimality of
all carrier denominators is unnecessary. This allows us to delete the
artificial scale $\mu=q^2+1$ and reuse the existing square $Y^2$.
The quotient substitution $Y=q^3s$, Pell doubling, and the final
product-square equation have predecessors in the pinned repository.

\subsection{Literal polynomial and fixed coefficients}

The only unknowns are
\[
 J,F,\alpha,i,T,w,s,n,\eta,\zeta\in\Z_{>0}.
\]
The fixed integers $\beta,K_0,M_C,M_F$ are the complete directly indexed
BFS943 tuple of package 133. They satisfy
\[
 \beta\ge31\text{ odd},\qquad K_0>0,\qquad
 0<M_C<\beta<M_F<2\beta.
\]
The supplied compressed unsigned little-endian payloads fix every byte;
\code{compact/rebuild_coefficients.py} reconstructs all four integers
from the unchanged finite tile specification and its 436-character code.
All four reconstructed raw hashes agree with the pinned parent.
The largest fixed port has 592,884 bits. This is the size of the stored
constant, not the height of the fully expanded polynomial.

All following symbols are polynomial abbreviations:
\begin{gather*}
 q=\beta J+1,\quad h=q-1,\quad C=q-F-\alpha,\quad X=wq,\\
 Y=q^3s,\quad k=\eta+\zeta,\quad c=kY+\eta,\quad D=Y^2,\\
 A=Y(X+1)+1,\quad \Delta=A^2-1,\quad H=4A-5,\\
 V_0=XD,\qquad E=1+V_0(V_0+1)k^2,\\
 R=[h(q-F)+\alpha](q^2-1)+(M_C+qM_F)J,\\
 e=k-R-1,\quad B=(K_0+w)C-F,\quad U=ic^2,\\
 a=\Delta U^2E,\quad F_0=1+a,\quad G=EF_0,\\
 F_*=1+4aF_0,\qquad Q=\Delta(F_*-1),\\
 b=c(T+\alpha)-R,\quad V=F_*b-c,\quad K_+=X(3c-2X)+2.
\end{gather*}
Choose $L=V^2$ for the 78-operation equation, or $L=bV$ for the
79-operation equation, and define
\begin{align*}
 M&=hDH,\\
 Z&=hD\bigl[H(n-LB)+K_+\bigr]-H(DB+he),\\
 R_G&=M^2G.
\end{align*}
The single equation is
\begin{equation}
 \boxed{\ P_L=Q(Z^2-R_GV^2)+(R_G-Z^2)=0.\ }
 \label{eq:compact-unscaled}
\end{equation}
There are no additional quantified variables, nonpolynomial primitives,
or implicit conditions in this definition. The actual circuits compute
$A-1$, $\Delta=(A-1)((A-1)+2)$, and $H=4(A-1)-1$.

\begin{theorem}[Unscaled compact consistency equations]
For either stated choice of $L$, ordinary arithmetic formalizable in ZFC,
composed with the pinned indexed native/refutation calibration, proves
\[
 \ZFC\vdash
 \bigl(\neg\exists\mathbf x\in\Z_{>0}^{10}\;
 P_L(\mathbf x)=0\bigr)\ \longleftrightarrow\ \Con(\ZFC).
\]
The exact degrees are 630 and 468, and the supplied evaluators have 78
and 79 live binary operations, respectively.
\end{theorem}

\subsection{Inherited interfaces and their scope}

The proof uses the following explicitly retained ordinary arithmetic
interfaces. They are statements of earlier work, not claims made by a
finite coefficient checker.
\begin{enumerate}
\item The universal main/auxiliary Pell folding and rank theorems:
from the strong norm, $c^2\mid I$, and the fixed-minus auxiliary
residues, they first recover $c=\psi_A(p)$ for an odd positive $p$,
and, once $0<p,R<c/2$, recover $p=R$.
\item The universal projection and fixed-minus parity theorems:
after the required index and size bounds, they give the whole-product
projection modulo $4A-5$ and $R\equiv3\pmod4$.
\item The zero-offset ratio estimate, proved uniformly for every
$\lambda\ge1$, and its half-binomial rounding consequences.
\item The indexed native coefficient calibration and fresh source
reconstruction: the resulting outer arithmetic data are equivalent to
a finite syntactic ZFC refutation. This branch retains its native/K215
ancestry; it is separate from the direct-calculus quartics.
\end{enumerate}
The exact source files and immutable commit are recorded in
\code{compact/SOURCE_PINS.json}. The new argument below verifies the
changed premises and supplies both existential directions. It does not
claim a new end-to-end proof-assistant derivation.

\subsection{Root extraction before index or sign semantics}

For proof notation only put
\[
 \delta=Q-1,\quad D_y=QV^2-1,\quad z=Z/M,\quad
 \kappa=K_+/H,\quad \rho=e/D,\quad t=B/h,\quad W=1+hL.
\]
The identities
\begin{equation}
 P_L=\delta Z^2-M^2GD_y,\qquad
 z=n+\kappa-Wt-\rho
 \label{eq:compact-carrier}
\end{equation}
hold without divisibility assumptions. Positivity of the inputs gives
$q\ge32$, $X\ge q$, $Y\ge q^3$, $k\ge2$, $c\ge3$, $A>4$,
$\Delta\ge3$, $U\ge c^2$, $E>1$, $F_*>c+1$, $Q>1$, and $M>0$.
For every integer $b$, $V=F_*b-c\ne0$: at $b=0$ it is $-c$;
otherwise $|V|\ge F_*-c>1$. Therefore $D_y>0$.

Identically,
\[
 F_*=(1+2a)^2,\qquad Q=4\Delta^2U^2G.
\]
Thus $G$ is coprime to both $\delta$ and $D_y$. A zero of
\eqref{eq:compact-unscaled} gives $\delta z^2=GD_y$.
For each prime dividing $G$, its valuation in $G$ is even, because
the other two factors are units at that prime and a rational square
has even valuation. Hence $G=g^2$ with $g>0$ integral. Since
$F_0\equiv1\pmod E$, the positive coprime factors split:
\[
 E=\tau^2,\qquad F_0=f_0^2,\qquad g=\tau f_0.
\]
Then $y=|z|/g$ is a positive rational satisfying
\begin{equation}
 QV^2-(Q-1)y^2=1.
 \label{eq:compact-adjacent}
\end{equation}

For completeness, the special adjacent form forces $y$ integral.
Writing $y=r/m$ in lowest terms gives $m^2\mid Q-1$. If
$v=|V|>1$, then $v<y<Qv/(Q-1)$. The norm-preserving transformation
\[
 v'=(2Q-1)v-2(Q-1)y,\qquad
 y'=(2Q-1)y-2Qv
\]
has an integer first coordinate with $0<|v'|<v$ and preserves the
reduced denominator $m$ of its second coordinate. Reflecting signs and
descending reaches $v=1,y=1$, so $m=1$. Consequently
\[
 z=\varepsilon gy\in\Z,\qquad \varepsilon\in\{-1,1\}.
\]
Set $f=1+2a$ and $I=2Ug$. Pell doubling gives
\begin{equation}
 f^2-\Delta I^2=1,\qquad Q=(\Delta I)^2,\qquad c^2\mid I,
 \qquad V=f^2[c(T+\alpha)-R]-c.
 \label{eq:compact-strong}
\end{equation}
No residue involving a scale $\mu$ has been used.

\subsection{Main folding and transport signs}

The folding interface applies to \eqref{eq:compact-adjacent} and
\eqref{eq:compact-strong}, since $c\le I$ and $V\equiv-c\pmod f$.
It gives $c=\psi_A(p)$ for odd $p\ge1$. The first norm yields
$k=2\psi_{2V_0+1}(n_0)$ for $n_0\ge1$. If $n_0=1$, then
$k=2$, $\eta=\zeta=1$, and $c=2Y+1$ lies strictly between the first
two main ordinates $1$ and $2A$, impossible. Therefore
\begin{equation}
 k\ge8V_0+4,\qquad
 c>kY\ge k>8XY^2\ge8q^7.
 \label{eq:compact-large-c}
\end{equation}
This replaces the earlier scale-based $8q^5$ bound by a stronger one.
It proves $c>X$ and $K_+>0$.
The unconditional mask estimate
$R<q^4+(\alpha+2)q^2$, together with
\eqref{eq:compact-large-c}, gives $b>c$ before any sign is assigned
to $B$. Thus $V>0$, $\kappa>0$, and
\begin{equation}
 b-e=c(T+\alpha)-k+1>0,\qquad \rho<b/D.
 \label{eq:compact-rho}
\end{equation}

For $L=V^2$, the identity
\[
 a^2-g^2=E\bigl[\Delta U^2E(\Delta U^2-1)-1\bigr]>0
\]
gives $g<a$, whence $V>(F_*-1)b>ab>3g$. Since
$y<3V/2$ by \eqref{eq:compact-adjacent}, $gy<L/2$.

For $L=bV$, the same fifth-index argument as in package 133 gives
$b>4g$. In detail, $c\ge3$ makes $p\ge3$, so $c\ge4\Delta+3$ and
$V>4Q$. Put $S=\Delta I$, so $Q=S^2\ge4$. Classification of
\eqref{eq:compact-adjacent} gives $SV=\chi_S(j)$ with odd $j$.
The first two odd indices give $V=1,4Q-3$, both excluded. Hence
\[
 V\ge16Q^2-20Q+5>8Q^2,
 \qquad b>4\Delta^2(F_*-1)=4\Delta^3I^2>4g.
\]
It follows that $gy<3bV/8<L/2$ here as well.

If $B\le-1$, then $Wt\le-L-1/h$. Equations
\eqref{eq:compact-carrier} and \eqref{eq:compact-rho} imply
\[
 n=\varepsilon gy-\kappa+Wt+\rho
   <L/2-L+b/D<0,
\]
because $V>b>1$ and $D>2$. This contradicts $n>0$. Therefore
$B\ge0$ and $C>0$. The outer mask bounds now give
$3q<R<q^4-2$. The rank interface has all its premises:
$c^2\mid I$, $c>2R$, $c>2p$, $V\equiv-c\pmod f$, and
$V\equiv-R\pmod c$. It gives $p=R$; the unchanged mask parity
then yields even $q$ and odd $J$.

\subsection{The new staged index lemma}

\begin{lemma}[Partial divisibility recovers the full denominator]
At every positive zero of \eqref{eq:compact-unscaled},
\[
 n_0=(R+1)/2\qquad\text{and}\qquad Y^2\mid e.
\]
Neither $\gcd(Y,q-1)=1$ nor integrality of $e/Y^2$ is needed as a premise.
\end{lemma}

\begin{proof}
The integral carrier in \eqref{eq:compact-carrier}, after clearing
$hH$, implies $Y^2\mid hHe$. Since $Y=q^3s$, one has $q^6\mid Y^2$.
Moreover $h=q-1$ and $H=4Y(X+1)-1\equiv-1\pmod q$, so
$\gcd(q^6,hH)=1$. Consequently $q^6\mid e$.

Let $P_0=2XY^2+1$. The base-one Pell recurrence gives
\[
 \psi_{P_0}(n_0)\equiv n_0\pmod{P_0-1},\qquad
 k\equiv2n_0\pmod{4XY^2}.
\]
It follows that $q^6\mid2n_0-R-1$.
But $P_0-A=Y(2XY-X-1)>0$. If $n_0\ge R$, monotonicity in the
base and index gives $k/2=\psi_{P_0}(n_0)\ge\psi_A(R)=c$,
contradicting $c=kY+\eta>k$. Thus $1\le n_0<R$, and
\[
 |2n_0-R-1|<R<q^4<q^6.
\]
The divisible integer must be zero. Returning to the \emph{full}
first-Pell congruence now gives $4XY^2\mid k-R-1=e$, and in
particular $Y^2\mid e$.
\end{proof}

The modulus $q^6$ is used only in this proof. It introduces no new
polynomial register. The whole denominator is recovered after the
index, which is the step that removes the artificial scale.

\subsection{Projection, source recovery, and positive completeness}

Write $R=2r+1$, with $r\ge48$. The retained zero-offset ratio lemma
holds for every $\lambda\ge1$ with
$A_\lambda=\lambda Y(X+1)+1$ and
$P_\lambda=2X\lambda^2Y^2+1$. At $\lambda=1$ it gives
\[
 \xi<\frac{\psi_A(2r+1)}{\psi_{P_0}(r+1)}
 <\xi\left(1+\frac1{Y(X+1)}\right)^{2r},\qquad
 \xi=\frac{(X+1)^{2r}}{X^r}.
\]
The lemma explicitly includes this value of $\lambda$. Its lower
comparison follows from
\[
 (2A_\lambda-1)^2X-(X+1)^2(2P_\lambda)
 =4\lambda Y(X+1)X+X-2(X+1)^2>0;
\]
the upper comparison uses
$A_\lambda\le\lambda(Y(X+1)+1)$.
Together with $Y<c/k<Y+1$, it yields
$Y>X^r/3$ and $X^{r+1}<3Y(X+1)<H$.

Now $\rho$ is integral, so the carrier gives $H\mid hK_+$.
The identity
\[
 K_+=HXc-2K_{\mathrm{quad}},\qquad
 K_{\mathrm{quad}}=X[X+2(A-2)c]-1
\]
allows cancellation of two through odd $H$. The universal projection
therefore implies
\[
 H\mid h(X-2^R)(2^RX+1).
\]
The retained whole-product bound applies with the same $X\ge q\ge32$
and $r\ge48$; its old multiplier $q^3h$ was larger. Thus the entire
product has absolute value below $X^{r+1}<H$, and $X=2^R$.
No cancellation of a possibly nonunit factor modulo a composite $H$
is assumed. It follows that $\kappa$ is integral. The remaining carrier
gives $h\mid WB$, hence $h\mid B$ because $W\equiv1\pmod h$.
The quotient $t\ge0$ is integral and $q^3\mid Y$ is literal.
Fixed-minus parity and exact half-binomial rounding now reconstruct
the earlier source with its \emph{fresh} base $Y(X+1)+2$.
The pinned native calibration yields a finite ZFC refutation.

Conversely, a refutation supplies canonical outer data with
$X=2^R$, $R\equiv3\pmod4$, $r\ge48$, $q^3\mid Y$, and
$t=B/h\ge0$, together with the half-binomial identity. Set $s=Y/q^3>0$.
At the new unscaled bases take
\[
 c=\psi_A(R),\quad k=2\psi_{P_0}(r+1),\quad
 \tau=\chi_{P_0}(r+1),\quad
 \eta=c-kY,\quad\zeta=k-\eta.
\]
The uniform ratio bound and the same rounding estimates prove
$0<\eta<k$, so both new coordinates are positive.
The base-one recurrence gives $e$ a positive multiple of $4XY^2$,
so $\rho\ge1$; projection gives a positive integral $\kappa$.

Choose a seed Pell pair $(f_0,I_0)$ at a positive multiple of the
finite matrix period modulo $c^2\tau$, and set
$i=I_0/(c^2\tau)>0$. Then $U=I_0/\tau$, $F_0=f_0^2$, and
$G=(\tau f_0)^2$. Double this pair:
\[
 f=2f_0^2-1,\qquad I=2f_0I_0=2Ug.
\]
The retained auxiliary-completeness theorem accepts this strong
pair without a restriction on the parity of its index. It supplies
$T'=T+\alpha>\alpha$ and positive $V,y$ satisfying the stated auxiliary
equation, so $T>0$. Finally choose
\[
 n=gy+Wt+\rho-\kappa.
\]
Every term is integral. Since
$f_0>(A-1)I_0>Xc>\kappa$, one has $g>\kappa$; hence $n>0$.
The resulting carrier is exactly $z=gy$, and the polynomial vanishes.
This constructs all ten positive inputs freshly and proves both
existential directions.

\subsection{Exact cost, degree, and checks}

The actual circuits have the following operation counts:
\begin{center}
\begin{tabular}{lrrrrr}
\toprule
Weight&Multiplications&Additions&Subtractions&Total&Degree\\
\midrule
$V^2$&46&18&14&78&630\\
$bV$&47&18&14&79&468\\
\bottomrule
\end{tabular}
\end{center}
The relevant global degrees are
\[
 \begin{array}{c|rrrrrrrrr}
 \text{register}&A&\Delta&E&U&a&F_*&Q&G&V\\\hline
 \deg&6&12&22&11&56&112&124&78&118
 \end{array}
\]
with $\deg b=6$, $\deg M=15$, and $\deg B=2$.
For $L=V^2$, $\deg Z=253$ and $QZ^2$ uniquely supplies degree 630.
For $L=bV$, $\deg Z=141$ and $-QM^2GV^2$ uniquely supplies degree 468.

Let $q_0=\beta J$, $k_0=\eta+\zeta$, and $t_0=T+\alpha$.
The entire leading homogeneous components are
\begin{align*}
 [P_{V^2}]_{630}
 &=2^{14}q_0^{382}w^{46}s^{112}i^{20}k_0^{64}t_0^4
       (q_0-F-\alpha)^2,\\
 [P_{bV}]_{468}
 &=-2^{10}q_0^{288}w^{34}s^{84}i^{14}k_0^{46}t_0^2.
\end{align*}
They are nonzero at the actual fixed positive $\beta$.
In particular the displayed formulas contain coefficients
$2^{14}\beta^{382}$ and $-2^{10}\beta^{288}$ on the respective
monomials
\[
 J^{382}w^{46}s^{112}i^{20}\eta^{64}T^4F^2,
 \qquad J^{288}w^{34}s^{84}i^{14}\eta^{46}T^2.
\]
These also illustrate the cost of the compact representation: the
expanded coefficients can be much larger than the stored ports.

The independent standard-library checker proves 28 exact polynomial
cut identities per circuit, covering all 78 or 79 actual registers.
It verifies topology, backward liveness, every fixed coefficient byte,
the complete leading forms (1,950 and 141 terms), and eight actual
algebraic corruptions per circuit. Normal and optimized Python checks
and byte-identical circuit rebuilds are supplied separately.
These finite checks verify the artifacts and their algebra; the
unbounded zero-existence equivalence is the ordinary proof above and
its stated inherited interfaces.


\section{What the verification proves}\label{sec:verification}

There are three distinct levels of evidence.

\textbf{Ordinary unbounded proofs.}
The new Pell, valuation, factorial, CRT, and cap arguments establish
the quantified mathematical implications. They compose with named
inherited arithmetic and logical interfaces. A finite test is not used
to infer a statement about all polynomial zeros.

\textbf{Exact finite algebra.}
The supplied independent checkers reconstruct every coefficient of
both quartics, verify every retained affine or monic elimination and
its restoration, check the integral translation in the signed branch,
and verify the compact circuits' topology, liveness, coefficient
binding, operation counts, and complete leading homogeneous forms.
The natural exporter independently checks 626 defining gates,
400 nonnegative affine restorations, and all 288 normalized rows before
the 106 further contractions. Its two nontrivial positivity restorations
are verified by exact rational sum-of-squares identities.

The natural reconstruction is replayed from the included relocated
arithmetic engine and its literal numerical-calculus input. The signed
replay starts from the included, pinned factored source of package 131;
it verifies the new source changes, projection, and final translation.
That starting boundary is deliberate and documented: it is not presented
as a rerun of every historical compiler stage. The compact coefficients
are reconstructed from their finite tile recipe, rather than accepted
from hashes alone. Checks run with explicit exceptions, and both normal
and optimized Python are exercised. Deliberately corrupted coefficients
or circuit data are rejected, including a natural coefficient alteration
whose claimed hash has been updated.

\textbf{Logical/formal boundary.}
We inspected and retained the concrete refutation interfaces and the
ordinary compiler proofs on which the new equations depend. We did not
rebuild the full transitive Lean development, generate a Rocq proof, or
produce an encoded object-language ZFC derivation of the new equations.
The result is an ordinary mathematical construction formalizable in ZFC,
with exact machine-checked finite algebra; it is not a newly completed
end-to-end proof-assistant development. The inherited interface names do
not turn the accompanying Python programs into a logical kernel.

\subsection{Reproduction}

From the extracted bundle root, using Python 3.10 or later:
\begin{quote}\ttfamily
python3 -B verify.py\par
python3 -B -O verify.py
\end{quote}
The root runner records the checks for each component. Individual
component READMEs state their exact replay boundaries and commands.
The ordinary mathematical arguments remain readable independently of
those programs. No search for a contradiction or a polynomial solution
is performed.

The complete paper can be compiled from \code{report/report.tex} with
pdfLaTeX twice. The final PDF has been rendered and visually inspected;
the archive also includes the full editable TeX sources. File hashes and
the immutable repository comparison commit are in the provenance files.
The data encoding for sparse polynomials is
\code{[coefficient, [coordinate indices]]}; repeated indices are repeated
factors, and an empty index list denotes a constant. All coordinate
indices start at zero. Thus the supplied files specify literal ordinary
polynomials without factorials, variable exponents, floor functions,
bounded-quantifier macros, or unspecified universal polynomials.

\section{Conclusions and the remaining target}

The strongest result here in the direct small-coefficient integer branch
is the 314-variable quartic. The new polynomial comparison center is
reusable wherever an indexed-Pell certificate already contains its
auxiliary norm. The natural branch gives a separate useful improvement
with a particularly short new factorial proof. The compact branch shows
that a denominator introduced to make a proof convenient can sometimes
be removed by changing the order in which arithmetic information is
recovered.

None of these observations supplies a lower bound establishing the
globally simplest consistency equation. The direct and compact
constructions still make different trades between number of unknowns,
degree, coefficient size, expanded support, and evaluation length.
Further useful work should preserve all these measures, together with
the domain and the exact calibration, when comparing a candidate.
The included literal equations and proof boundaries make such a
comparison reproducible.

\begin{thebibliography}{9}
\bibitem{repo} V. Reshetnikov, \emph{ProveIt}, comparison commit
\code{1c5367452a0bd4dd7e9e4bceb1508b62dbd2ff61}.
\url{https://github.com/VladimirReshetnikov/ProveIt}.
\bibitem{parent129} ProveIt package 129, \emph{Semantic Anchors for
Diophantine Consistency}, 24 October 2026, especially the natural
233/234 equations and ordinary arithmetic proof.
\url{https://github.com/VladimirReshetnikov/ProveIt/tree/1c5367452a0bd4dd7e9e4bceb1508b62dbd2ff61/Computability/DiophantineConsistency/con-zfc-diophantine-forms/129-semantic-pell-anchor}.
\bibitem{parent131} ProveIt package 131, \emph{Shared Pell Caps and
Coprime Product Equations for ZFC Consistency}, signed 324-variable
endpoint and its retained calibration, same pinned commit.
\bibitem{parent132} ProveIt package 132, \emph{Signed Pell Caps and
Free Scale Carriers}, odd-quotient and signed-cap arguments,
same pinned commit.
\bibitem{parent133} ProveIt package 133, \emph{Staged Pell Certificates
and a Product-Square Equation for Con(ZFC)}, compact 84/436 and
85/328 endpoints and complete coefficient tuple, same pinned commit.
\bibitem{cap128} ProveIt package 128, \emph{One Quadratic Constraint
for Many Pell Certificates}, one-row cap theorem, same pinned commit.
\bibitem{jones} J. P. Jones, Universal Diophantine equation,
\emph{Journal of Symbolic Logic} \textbf{47} (1982), 549--571.
\bibitem{sun} Z.-W. Sun, \emph{Ten unknowns for Hilbert's tenth
problem over the integers}, arXiv:2609.01594v2, 2 September 2026.
\url{https://arxiv.org/abs/2609.01594}.
\end{thebibliography}

\appendix
\section{The complete unrestricted-integer residual list}
\label{app:integer}
All 314 variables \(z_0,\ldots,z_{313}\) range independently over
\(\Z\). The equation is \(\sum_{j=0}^{203}r_j^2=0\), using every
row below. No semantic abbreviation remains in this finite definition.
% Generated from direct/equation.json; all residuals are printed.
\begingroup
\small
\setlength{\tabcolsep}{2pt}
\begin{longtable}{@{}r@{\;}l@{}}
\(r_{0}=\) & \(\displaystyle -z_{3}^{2} +z_{4}\) \\[3pt]
\(r_{1}=\) & \(\displaystyle -1 -z_{0}z_{2} +z_{9}\) \\[3pt]
\(r_{2}=\) & \(\displaystyle -1 +z_{1} +z_{2}\) \\
 & \(\displaystyle -z_{2}z_{7} -z_{3} -z_{7}\) \\[3pt]
\(r_{3}=\) & \(\displaystyle -2 +z_{1} -z_{8}z_{9}\) \\[3pt]
\(r_{4}=\) & \(\displaystyle 1 +z_{4} -z_{5}z_{20}\) \\
 & \(\displaystyle +z_{18}\) \\[3pt]
\(r_{5}=\) & \(\displaystyle -1 +2z_{10} -z_{10}^{2}\) \\
 & \(\displaystyle -2z_{10}z_{11} +2z_{11} -z_{11}^{2}\) \\
 & \(\displaystyle +z_{21}\) \\[3pt]
\(r_{6}=\) & \(\displaystyle -5 -30z_{10} -36z_{10}^{2}\) \\
 & \(\displaystyle -96z_{10}z_{11} -24z_{10}z_{21} -40z_{11}\) \\
 & \(\displaystyle -64z_{11}^{2} -32z_{11}z_{21} -10z_{21}\) \\
 & \(\displaystyle -4z_{21}^{2} +z_{22}\) \\[3pt]
\(r_{7}=\) & \(\displaystyle -2 +2z_{10} -z_{10}^{2}\) \\
 & \(\displaystyle -4z_{10}z_{22} +4z_{22} -4z_{22}^{2}\) \\
 & \(\displaystyle +z_{23}\) \\[3pt]
\(r_{8}=\) & \(\displaystyle 1 +6z_{10} -36z_{10}^{2}\) \\
 & \(\displaystyle -192z_{10}z_{22} -24z_{10}z_{23} +16z_{22}\) \\
 & \(\displaystyle -256z_{22}^{2} -64z_{22}z_{23} +2z_{23}\) \\
 & \(\displaystyle -4z_{23}^{2} +z_{24}\) \\[3pt]
\(r_{9}=\) & \(\displaystyle -2 -z_{11}^{2} -2z_{11}z_{12}\) \\
 & \(\displaystyle -z_{12}^{2} +z_{25}\) \\[3pt]
\(r_{10}=\) & \(\displaystyle -9 +7z_{10} -z_{10}^{2}\) \\
 & \(\displaystyle -4z_{10}z_{11} -4z_{10}z_{25} +12z_{11}\) \\
 & \(\displaystyle -4z_{11}^{2} -8z_{11}z_{25} +12z_{25}\) \\
 & \(\displaystyle -4z_{25}^{2} +z_{26}\) \\[3pt]
\(r_{11}=\) & \(\displaystyle -1 +2z_{10} -z_{10}^{2}\) \\
 & \(\displaystyle -2z_{10}z_{12} +2z_{12} -z_{12}^{2}\) \\
 & \(\displaystyle +z_{27}\) \\[3pt]
\(r_{12}=\) & \(\displaystyle -6z_{10} -144z_{10}^{2} -192z_{10}z_{11}\) \\
 & \(\displaystyle -192z_{10}z_{12} -48z_{10}z_{21} -48z_{10}z_{27}\) \\
 & \(\displaystyle -8z_{11} -64z_{11}^{2} -128z_{11}z_{12}\) \\
 & \(\displaystyle -32z_{11}z_{21} -32z_{11}z_{27} -64z_{12}^{2}\) \\
 & \(\displaystyle -32z_{12}z_{21} -32z_{12}z_{27} -2z_{21}\) \\
 & \(\displaystyle -4z_{21}^{2} -8z_{21}z_{27} -4z_{27}^{2}\) \\
 & \(\displaystyle +z_{28}\) \\[3pt]
\(r_{13}=\) & \(\displaystyle -4z_{26}^{2} -8z_{26}z_{28} -4z_{28}^{2}\) \\
 & \(\displaystyle +z_{29}\) \\[3pt]
\(r_{14}=\) & \(\displaystyle -6 +16z_{12} -16z_{12}^{2}\) \\
 & \(\displaystyle -16z_{12}z_{25} -16z_{12}z_{27} +8z_{25}\) \\
 & \(\displaystyle -4z_{25}^{2} -8z_{25}z_{27} +8z_{27}\) \\
 & \(\displaystyle -4z_{27}^{2} +z_{30}\) \\[3pt]
\(r_{15}=\) & \(\displaystyle 1 +6z_{10} -36z_{10}^{2}\) \\
 & \(\displaystyle -96z_{10}z_{11} -48z_{10}z_{12} -24z_{10}z_{21}\) \\
 & \(\displaystyle -48z_{10}z_{27} -24z_{10}z_{30} +8z_{11}\) \\
 & \(\displaystyle -64z_{11}^{2} -64z_{11}z_{12} -32z_{11}z_{21}\) \\
 & \(\displaystyle -64z_{11}z_{27} -32z_{11}z_{30} +8z_{12}\) \\
 & \(\displaystyle -16z_{12}^{2} -16z_{12}z_{21} -32z_{12}z_{27}\) \\
 & \(\displaystyle -16z_{12}z_{30} +2z_{21} -4z_{21}^{2}\) \\
 & \(\displaystyle -16z_{21}z_{27} -8z_{21}z_{30} +8z_{27}\) \\
 & \(\displaystyle -16z_{27}^{2} -16z_{27}z_{30} +4z_{30}\) \\
 & \(\displaystyle -4z_{30}^{2} +z_{31}\) \\[3pt]
\(r_{16}=\) & \(\displaystyle -10 -5z_{10} -z_{10}^{2}\) \\
 & \(\displaystyle -12z_{10}z_{11} -16z_{10}z_{12} -4z_{10}z_{25}\) \\
 & \(\displaystyle -36z_{11} -36z_{11}^{2} -96z_{11}z_{12}\) \\
 & \(\displaystyle -24z_{11}z_{25} -48z_{12} -64z_{12}^{2}\) \\
 & \(\displaystyle -32z_{12}z_{25} -12z_{25} -4z_{25}^{2}\) \\
 & \(\displaystyle +z_{32}\) \\[3pt]
\(r_{17}=\) & \(\displaystyle -2 -24z_{10} -36z_{10}^{2}\) \\
 & \(\displaystyle -12z_{10}z_{11} -96z_{10}z_{12} -24z_{10}z_{27}\) \\
 & \(\displaystyle -3z_{11} -z_{11}^{2} -16z_{11}z_{12}\) \\
 & \(\displaystyle -4z_{11}z_{27} -32z_{12} -64z_{12}^{2}\) \\
 & \(\displaystyle -32z_{12}z_{27} -8z_{27} -4z_{27}^{2}\) \\
 & \(\displaystyle +z_{33}\) \\[3pt]
\(r_{18}=\) & \(\displaystyle 2 -4z_{32}^{2} -8z_{32}z_{33}\) \\
 & \(\displaystyle -4z_{33}^{2} +z_{34}\) \\[3pt]
\(r_{19}=\) & \(\displaystyle -12 +16z_{10} -4z_{10}^{2}\) \\
 & \(\displaystyle -8z_{10}z_{12} -8z_{10}z_{21} -8z_{10}z_{25}\) \\
 & \(\displaystyle +16z_{12} -4z_{12}^{2} -8z_{12}z_{21}\) \\
 & \(\displaystyle -8z_{12}z_{25} +16z_{21} -4z_{21}^{2}\) \\
 & \(\displaystyle -8z_{21}z_{25} +16z_{25} -4z_{25}^{2}\) \\
 & \(\displaystyle +z_{35}\) \\[3pt]
\(r_{20}=\) & \(\displaystyle -3 -24z_{10} -64z_{10}^{2}\) \\
 & \(\displaystyle -160z_{10}z_{12} -64z_{10}z_{25} -32z_{10}z_{27}\) \\
 & \(\displaystyle -32z_{10}z_{35} -26z_{12} -100z_{12}^{2}\) \\
 & \(\displaystyle -80z_{12}z_{25} -40z_{12}z_{27} -40z_{12}z_{35}\) \\
 & \(\displaystyle -8z_{25} -16z_{25}^{2} -16z_{25}z_{27}\) \\
 & \(\displaystyle -16z_{25}z_{35} -6z_{27} -4z_{27}^{2}\) \\
 & \(\displaystyle -8z_{27}z_{35} -4z_{35} -4z_{35}^{2}\) \\
 & \(\displaystyle +z_{36}\) \\[3pt]
\(r_{21}=\) & \(\displaystyle -1 -18z_{10} -36z_{10}^{2}\) \\
 & \(\displaystyle -96z_{10}z_{11} -24z_{10}z_{21} -24z_{11}\) \\
 & \(\displaystyle -64z_{11}^{2} -32z_{11}z_{21} -6z_{21}\) \\
 & \(\displaystyle -4z_{21}^{2} +z_{37}\) \\[3pt]
\(r_{22}=\) & \(\displaystyle -2 +2z_{10} -z_{10}^{2}\) \\
 & \(\displaystyle -4z_{10}z_{37} +4z_{37} -4z_{37}^{2}\) \\
 & \(\displaystyle +z_{38}\) \\[3pt]
\(r_{23}=\) & \(\displaystyle 1 +6z_{10} -36z_{10}^{2}\) \\
 & \(\displaystyle -192z_{10}z_{37} -24z_{10}z_{38} +16z_{37}\) \\
 & \(\displaystyle -256z_{37}^{2} -64z_{37}z_{38} +2z_{38}\) \\
 & \(\displaystyle -4z_{38}^{2} +z_{39}\) \\[3pt]
\(r_{24}=\) & \(\displaystyle -z_{11}^{2} -4z_{11}z_{39} -4z_{39}^{2}\) \\
 & \(\displaystyle +z_{40}\) \\[3pt]
\(r_{25}=\) & \(\displaystyle -6 -30z_{10} -64z_{10}^{2}\) \\
 & \(\displaystyle -128z_{10}z_{11} -64z_{10}z_{21} -32z_{11}\) \\
 & \(\displaystyle -64z_{11}^{2} -64z_{11}z_{21} -14z_{21}\) \\
 & \(\displaystyle -16z_{21}^{2} +z_{41}\) \\[3pt]
\(r_{26}=\) & \(\displaystyle -4 +8z_{11} -64z_{11}^{2}\) \\
 & \(\displaystyle -96z_{11}z_{12} -32z_{11}z_{25} -32z_{11}z_{41}\) \\
 & \(\displaystyle +6z_{12} -36z_{12}^{2} -24z_{12}z_{25}\) \\
 & \(\displaystyle -24z_{12}z_{41} +2z_{25} -4z_{25}^{2}\) \\
 & \(\displaystyle -8z_{25}z_{41} +4z_{41} -4z_{41}^{2}\) \\
 & \(\displaystyle +z_{42}\) \\[3pt]
\(r_{27}=\) & \(\displaystyle 8z_{10} -64z_{10}^{2} -96z_{10}z_{12}\) \\
 & \(\displaystyle -32z_{10}z_{27} -32z_{10}z_{42} +6z_{12}\) \\
 & \(\displaystyle -36z_{12}^{2} -24z_{12}z_{27} -24z_{12}z_{42}\) \\
 & \(\displaystyle +2z_{27} -4z_{27}^{2} -8z_{27}z_{42}\) \\
 & \(\displaystyle +4z_{42} -4z_{42}^{2} +z_{43}\) \\[3pt]
\(r_{28}=\) & \(\displaystyle -z_{12}^{2} -2z_{12}z_{13} -z_{13}^{2}\) \\
 & \(\displaystyle +z_{44}\) \\[3pt]
\(r_{29}=\) & \(\displaystyle 2z_{10} -z_{10}^{2} -4z_{10}z_{12}\) \\
 & \(\displaystyle -4z_{10}z_{44} +4z_{12} -4z_{12}^{2}\) \\
 & \(\displaystyle -8z_{12}z_{44} +4z_{44} -4z_{44}^{2}\) \\
 & \(\displaystyle +z_{45}\) \\[3pt]
\(r_{30}=\) & \(\displaystyle -35 +24z_{10} -4z_{10}^{2}\) \\
 & \(\displaystyle -8z_{10}z_{11} -8z_{10}z_{25} -8z_{10}z_{27}\) \\
 & \(\displaystyle +24z_{11} -4z_{11}^{2} -8z_{11}z_{25}\) \\
 & \(\displaystyle -8z_{11}z_{27} +24z_{25} -4z_{25}^{2}\) \\
 & \(\displaystyle -8z_{25}z_{27} +24z_{27} -4z_{27}^{2}\) \\
 & \(\displaystyle +z_{46}\) \\[3pt]
\(r_{31}=\) & \(\displaystyle -16z_{10}^{2} -32z_{10}z_{12} -32z_{10}z_{27}\) \\
 & \(\displaystyle -32z_{10}z_{44} -16z_{10}z_{45} -16z_{10}z_{46}\) \\
 & \(\displaystyle -4z_{12} -16z_{12}^{2} -32z_{12}z_{27}\) \\
 & \(\displaystyle -32z_{12}z_{44} -16z_{12}z_{45} -16z_{12}z_{46}\) \\
 & \(\displaystyle -16z_{27}^{2} -32z_{27}z_{44} -16z_{27}z_{45}\) \\
 & \(\displaystyle -16z_{27}z_{46} -4z_{44} -16z_{44}^{2}\) \\
 & \(\displaystyle -16z_{44}z_{45} -16z_{44}z_{46} -2z_{45}\) \\
 & \(\displaystyle -4z_{45}^{2} -8z_{45}z_{46} -4z_{46}^{2}\) \\
 & \(\displaystyle +z_{47}\) \\[3pt]
\(r_{32}=\) & \(\displaystyle z_{12} -25z_{12}^{2} -40z_{12}z_{27}\) \\
 & \(\displaystyle -20z_{12}z_{30} -16z_{27}^{2} -16z_{27}z_{30}\) \\
 & \(\displaystyle -4z_{30}^{2} +z_{48}\) \\[3pt]
\(r_{33}=\) & \(\displaystyle -1 +8z_{11} -4z_{11}^{2}\) \\
 & \(\displaystyle -8z_{11}z_{25} +8z_{25} -4z_{25}^{2}\) \\
 & \(\displaystyle +z_{49}\) \\[3pt]
\(r_{34}=\) & \(\displaystyle -1 -16z_{11}^{2} -8z_{11}z_{12}\) \\
 & \(\displaystyle -32z_{11}z_{25} -16z_{11}z_{49} -z_{12}^{2}\) \\
 & \(\displaystyle -8z_{12}z_{25} -4z_{12}z_{49} -16z_{25}^{2}\) \\
 & \(\displaystyle -16z_{25}z_{49} -4z_{49}^{2} +z_{50}\) \\[3pt]
\(r_{35}=\) & \(\displaystyle 1 +z_{11} -z_{11}^{2}\) \\
 & \(\displaystyle -4z_{11}z_{12} -4z_{11}z_{50} -4z_{12}^{2}\) \\
 & \(\displaystyle -8z_{12}z_{50} -4z_{50}^{2} +z_{51}\) \\[3pt]
\(r_{36}=\) & \(\displaystyle -1 -4z_{51} -4z_{51}^{2}\) \\
 & \(\displaystyle +z_{52}\) \\[3pt]
\(r_{37}=\) & \(\displaystyle -1 -z_{12}^{2} -8z_{12}z_{51}\) \\
 & \(\displaystyle -4z_{12}z_{52} -16z_{51}^{2} -16z_{51}z_{52}\) \\
 & \(\displaystyle -4z_{52}^{2} +z_{53}\) \\[3pt]
\(r_{38}=\) & \(\displaystyle 2 +z_{11} -z_{11}^{2}\) \\
 & \(\displaystyle -4z_{11}z_{12} -4z_{11}z_{53} -4z_{12}^{2}\) \\
 & \(\displaystyle -8z_{12}z_{53} -4z_{53}^{2} +z_{54}\) \\[3pt]
\(r_{39}=\) & \(\displaystyle -4z_{48}^{2} -8z_{48}z_{54} -4z_{54}^{2}\) \\
 & \(\displaystyle +z_{55}\) \\[3pt]
\(r_{40}=\) & \(\displaystyle -9 +7z_{10} -z_{10}^{2}\) \\
 & \(\displaystyle -8z_{10}z_{48} -4z_{10}z_{55} +24z_{48}\) \\
 & \(\displaystyle -16z_{48}^{2} -16z_{48}z_{55} +12z_{55}\) \\
 & \(\displaystyle -4z_{55}^{2} +z_{56}\) \\[3pt]
\(r_{41}=\) & \(\displaystyle -1 -4z_{56} -4z_{56}^{2}\) \\
 & \(\displaystyle +z_{57}\) \\[3pt]
\(r_{42}=\) & \(\displaystyle -z_{12}^{2} -8z_{12}z_{56} -4z_{12}z_{57}\) \\
 & \(\displaystyle -16z_{56}^{2} -16z_{56}z_{57} -4z_{57}^{2}\) \\
 & \(\displaystyle +z_{58}\) \\[3pt]
\(r_{43}=\) & \(\displaystyle -z_{10} -z_{10}^{2} -4z_{10}z_{12}\) \\
 & \(\displaystyle -4z_{10}z_{58} -4z_{12} -4z_{12}^{2}\) \\
 & \(\displaystyle -8z_{12}z_{58} -4z_{58} -4z_{58}^{2}\) \\
 & \(\displaystyle +z_{59}\) \\[3pt]
\(r_{44}=\) & \(\displaystyle -10 -12z_{59} -4z_{59}^{2}\) \\
 & \(\displaystyle +z_{60}\) \\[3pt]
\(r_{45}=\) & \(\displaystyle 1 +4z_{11} -4z_{11}^{2}\) \\
 & \(\displaystyle -8z_{11}z_{25} -16z_{11}z_{59} -8z_{11}z_{60}\) \\
 & \(\displaystyle +4z_{25} -4z_{25}^{2} -16z_{25}z_{59}\) \\
 & \(\displaystyle -8z_{25}z_{60} +8z_{59} -16z_{59}^{2}\) \\
 & \(\displaystyle -16z_{59}z_{60} +4z_{60} -4z_{60}^{2}\) \\
 & \(\displaystyle +z_{61}\) \\[3pt]
\(r_{46}=\) & \(\displaystyle z_{11} -25z_{11}^{2} -40z_{11}z_{25}\) \\
 & \(\displaystyle -20z_{11}z_{61} -16z_{25}^{2} -16z_{25}z_{61}\) \\
 & \(\displaystyle -4z_{61}^{2} +z_{62}\) \\[3pt]
\(r_{47}=\) & \(\displaystyle -10 -12z_{54} -4z_{54}^{2}\) \\
 & \(\displaystyle +z_{63}\) \\[3pt]
\(r_{48}=\) & \(\displaystyle 1 -16z_{54}^{2} -16z_{54}z_{62}\) \\
 & \(\displaystyle -16z_{54}z_{63} -4z_{62}^{2} -8z_{62}z_{63}\) \\
 & \(\displaystyle -4z_{63}^{2} +z_{64}\) \\[3pt]
\(r_{49}=\) & \(\displaystyle z_{12} -z_{12}^{2} -8z_{12}z_{62}\) \\
 & \(\displaystyle -4z_{12}z_{64} -16z_{62}^{2} -16z_{62}z_{64}\) \\
 & \(\displaystyle -4z_{64}^{2} +z_{65}\) \\[3pt]
\(r_{50}=\) & \(\displaystyle -4z_{10} -4z_{10}^{2} -4z_{10}z_{12}\) \\
 & \(\displaystyle -8z_{10}z_{21} -z_{12} -z_{12}^{2}\) \\
 & \(\displaystyle -4z_{12}z_{21} -4z_{21} -4z_{21}^{2}\) \\
 & \(\displaystyle +z_{66}\) \\[3pt]
\(r_{51}=\) & \(\displaystyle 1 +4z_{10} -4z_{10}^{2}\) \\
 & \(\displaystyle -8z_{10}z_{21} -8z_{10}z_{66} +4z_{21}\) \\
 & \(\displaystyle -4z_{21}^{2} -8z_{21}z_{66} +4z_{66}\) \\
 & \(\displaystyle -4z_{66}^{2} +z_{67}\) \\[3pt]
\(r_{52}=\) & \(\displaystyle -2 -4z_{10} -16z_{10}^{2}\) \\
 & \(\displaystyle -32z_{10}z_{21} -16z_{10}z_{65} -16z_{10}z_{67}\) \\
 & \(\displaystyle -4z_{21} -16z_{21}^{2} -16z_{21}z_{65}\) \\
 & \(\displaystyle -16z_{21}z_{67} -4z_{65}^{2} -8z_{65}z_{67}\) \\
 & \(\displaystyle -2z_{67} -4z_{67}^{2} +z_{68}\) \\[3pt]
\(r_{53}=\) & \(\displaystyle -2 -4z_{68} -4z_{68}^{2}\) \\
 & \(\displaystyle +z_{69}\) \\[3pt]
\(r_{54}=\) & \(\displaystyle -z_{3}z_{15} +z_{70}\) \\[3pt]
\(r_{55}=\) & \(\displaystyle z_{14} -z_{15} -2z_{24}z_{70}\) \\
 & \(\displaystyle -z_{71}\) \\[3pt]
\(r_{56}=\) & \(\displaystyle z_{14} -z_{15} -4z_{26}z_{70}\) \\
 & \(\displaystyle -2z_{29}z_{70} -z_{72}\) \\[3pt]
\(r_{57}=\) & \(\displaystyle z_{14} -z_{15} -2z_{31}z_{70}\) \\
 & \(\displaystyle -6z_{70} -z_{73}\) \\[3pt]
\(r_{58}=\) & \(\displaystyle z_{14} -z_{15} -4z_{32}z_{70}\) \\
 & \(\displaystyle -2z_{34}z_{70} -z_{74}\) \\[3pt]
\(r_{59}=\) & \(\displaystyle z_{14} -z_{15} -2z_{36}z_{70}\) \\
 & \(\displaystyle +6z_{70} -z_{75}\) \\[3pt]
\(r_{60}=\) & \(\displaystyle z_{14} -z_{15} -4z_{39}z_{70}\) \\
 & \(\displaystyle -2z_{40}z_{70} -z_{76}\) \\[3pt]
\(r_{61}=\) & \(\displaystyle z_{14} -z_{15} -2z_{43}z_{70}\) \\
 & \(\displaystyle -z_{77}\) \\[3pt]
\(r_{62}=\) & \(\displaystyle z_{14} -z_{15} -2z_{47}z_{70}\) \\
 & \(\displaystyle -4z_{70} -z_{78}\) \\[3pt]
\(r_{63}=\) & \(\displaystyle -z_{10}z_{11} +z_{10}z_{12} +z_{11}\) \\
 & \(\displaystyle +z_{11}z_{12} -z_{12} -z_{12}^{2}\) \\
 & \(\displaystyle +z_{79}\) \\[3pt]
\(r_{64}=\) & \(\displaystyle 1 +2z_{13} +z_{13}^{2}\) \\
 & \(\displaystyle -z_{79}^{2} +z_{80}\) \\[3pt]
\(r_{65}=\) & \(\displaystyle -z_{80}^{2} +z_{81}\) \\[3pt]
\(r_{66}=\) & \(\displaystyle -4z_{68}z_{70} -2z_{69}z_{70} +2z_{70}\) \\
 & \(\displaystyle +z_{82}\) \\[3pt]
\(r_{67}=\) & \(\displaystyle -z_{14}^{2} +2z_{14}z_{15} +2z_{14}z_{82}\) \\
 & \(\displaystyle -z_{15}^{2} -2z_{15}z_{82} -z_{82}^{2}\) \\
 & \(\displaystyle +z_{83}\) \\[3pt]
\(r_{68}=\) & \(\displaystyle -1 -z_{3}z_{11} +z_{84}\) \\[3pt]
\(r_{69}=\) & \(\displaystyle -z_{15}z_{84} +z_{86}z_{87} +z_{88}\) \\[3pt]
\(r_{70}=\) & \(\displaystyle -1 +z_{3} -z_{3}z_{10}\) \\
 & \(\displaystyle +z_{85}\) \\[3pt]
\(r_{71}=\) & \(\displaystyle -1 +2z_{3} -2z_{3}z_{10}\) \\
 & \(\displaystyle -2z_{3}z_{21} +z_{86}\) \\[3pt]
\(r_{72}=\) & \(\displaystyle -z_{14}z_{85} +z_{87}\) \\[3pt]
\(r_{73}=\) & \(\displaystyle -1 -6z_{3}z_{10} -8z_{3}z_{11}\) \\
 & \(\displaystyle -2z_{3}z_{21} +z_{89}\) \\[3pt]
\(r_{74}=\) & \(\displaystyle z_{14}z_{84} -z_{15}z_{89} +z_{90}\) \\[3pt]
\(r_{75}=\) & \(\displaystyle -z_{14}z_{71} +z_{15}z_{71} +z_{91}\) \\[3pt]
\(r_{76}=\) & \(\displaystyle -z_{72}z_{91} +z_{92}\) \\[3pt]
\(r_{77}=\) & \(\displaystyle -z_{73}z_{92} +z_{93}\) \\[3pt]
\(r_{78}=\) & \(\displaystyle -z_{74}z_{93} +z_{94}\) \\[3pt]
\(r_{79}=\) & \(\displaystyle -z_{75}z_{94} +z_{95}\) \\[3pt]
\(r_{80}=\) & \(\displaystyle -z_{76}z_{95} +z_{96}\) \\[3pt]
\(r_{81}=\) & \(\displaystyle -z_{77}z_{96} +z_{97}\) \\[3pt]
\(r_{82}=\) & \(\displaystyle -z_{78}z_{97} +z_{98}\) \\[3pt]
\(r_{83}=\) & \(\displaystyle -z_{88}z_{98} +z_{99}\) \\[3pt]
\(r_{84}=\) & \(\displaystyle -z_{90}z_{99} +z_{100}\) \\[3pt]
\(r_{85}=\) & \(\displaystyle -z_{100}^{2} +z_{101}\) \\[3pt]
\(r_{86}=\) & \(\displaystyle -1 -z_{2} -z_{2}z_{18}\) \\
 & \(\displaystyle +z_{102}\) \\[3pt]
\(r_{87}=\) & \(\displaystyle -1 -z_{16}z_{102} +z_{103}\) \\[3pt]
\(r_{88}=\) & \(\displaystyle -2z_{1}^{2} +2z_{1}z_{14} +2z_{1}z_{15}\) \\
 & \(\displaystyle +2z_{1}z_{103} +2z_{1}z_{104} -z_{14}^{2}\) \\
 & \(\displaystyle -2z_{14}z_{103} -z_{15}^{2} -2z_{15}z_{104}\) \\
 & \(\displaystyle +z_{19}z_{106} -z_{81}z_{101} -z_{83}z_{101}\) \\
 & \(\displaystyle -z_{103}^{2} -z_{104}^{2}\) \\[3pt]
\(r_{89}=\) & \(\displaystyle -1 -z_{2}z_{17} -z_{17}z_{102}\) \\
 & \(\displaystyle +z_{104}\) \\[3pt]
\(r_{90}=\) & \(\displaystyle -9z_{10} -9z_{11} -4z_{21}\) \\
 & \(\displaystyle -2z_{24} -4z_{26} -2z_{29}\) \\
 & \(\displaystyle -2z_{31} -4z_{32} -2z_{34}\) \\
 & \(\displaystyle -2z_{36} -4z_{39} -2z_{40}\) \\
 & \(\displaystyle -2z_{43} -2z_{47} -4z_{68}\) \\
 & \(\displaystyle -2z_{69} +z_{105}\) \\[3pt]
\(r_{91}=\) & \(\displaystyle -z_{107}z_{110} +z_{110} +z_{112}\) \\[3pt]
\(r_{92}=\) & \(\displaystyle -z_{5}z_{109} +z_{108} -z_{109}z_{112}\) \\
 & \(\displaystyle -z_{111}\) \\[3pt]
\(r_{93}=\) & \(\displaystyle -z_{19}z_{114} +z_{117}\) \\[3pt]
\(r_{94}=\) & \(\displaystyle z_{113} -z_{115}z_{117} -z_{116}\) \\[3pt]
\(r_{95}=\) & \(\displaystyle -z_{117}z_{119} +z_{118} -z_{120}\) \\[3pt]
\(r_{96}=\) & \(\displaystyle -z_{117}z_{122} +z_{121} -z_{123}\) \\[3pt]
\(r_{97}=\) & \(\displaystyle -z_{117}z_{125} +z_{124} -z_{126}\) \\[3pt]
\(r_{98}=\) & \(\displaystyle -z_{19}z_{128} +z_{131}\) \\[3pt]
\(r_{99}=\) & \(\displaystyle z_{127} -z_{129}z_{131} -z_{130}\) \\[3pt]
\(r_{100}=\) & \(\displaystyle -z_{131}z_{133} +z_{132} -z_{134}\) \\[3pt]
\(r_{101}=\) & \(\displaystyle -z_{19}z_{136} +z_{139}\) \\[3pt]
\(r_{102}=\) & \(\displaystyle z_{135} -z_{137}z_{139} -z_{138}\) \\[3pt]
\(r_{103}=\) & \(\displaystyle -z_{139}z_{141} +z_{140} -z_{142}\) \\[3pt]
\(r_{104}=\) & \(\displaystyle -z_{19}z_{109} +z_{146}\) \\[3pt]
\(r_{105}=\) & \(\displaystyle z_{143} -z_{144}z_{146} -z_{145}\) \\[3pt]
\(r_{106}=\) & \(\displaystyle -z_{117}z_{148} +z_{147} -z_{149}\) \\[3pt]
\(r_{107}=\) & \(\displaystyle z_{150} -2z_{150}z_{152} +z_{153}\) \\[3pt]
\(r_{108}=\) & \(\displaystyle -z_{150}z_{153} -z_{153} +z_{154}\) \\[3pt]
\(r_{109}=\) & \(\displaystyle 2z_{151} +z_{151}^{2} +z_{153}^{2}\) \\
 & \(\displaystyle -z_{154}^{2}\) \\[3pt]
\(r_{110}=\) & \(\displaystyle 1 +z_{151} -z_{151}z_{156}\) \\
 & \(\displaystyle -z_{156} +z_{157}\) \\[3pt]
\(r_{111}=\) & \(\displaystyle z_{155}^{2} -2z_{156} +z_{156}^{2}\) \\
 & \(\displaystyle -z_{157}^{2}\) \\[3pt]
\(r_{112}=\) & \(\displaystyle z_{4}^{2} -2z_{4}z_{151} -2z_{151}\) \\
 & \(\displaystyle +z_{158}\) \\[3pt]
\(r_{113}=\) & \(\displaystyle 1 +z_{4} -z_{4}z_{156}\) \\
 & \(\displaystyle +z_{6} -z_{155} -z_{156}\) \\
 & \(\displaystyle +z_{157} +z_{158}z_{159}\) \\[3pt]
\(r_{114}=\) & \(\displaystyle -z_{151}z_{161} -z_{161} +z_{162}\) \\[3pt]
\(r_{115}=\) & \(\displaystyle -1 +z_{160}^{2} +z_{161}^{2}\) \\
 & \(\displaystyle -z_{162}^{2}\) \\[3pt]
\(r_{116}=\) & \(\displaystyle 1 -2z_{5} +z_{5}^{2}\) \\
 & \(\displaystyle -2z_{5}z_{151} +z_{163}\) \\[3pt]
\(r_{117}=\) & \(\displaystyle -z_{5}z_{161} +z_{19} -z_{160}\) \\
 & \(\displaystyle +z_{162} +z_{163}z_{164}\) \\[3pt]
\(r_{118}=\) & \(\displaystyle -1 +2z_{151}z_{224} -z_{165}\) \\
 & \(\displaystyle +2z_{224} -z_{224}^{2}\) \\[3pt]
\(r_{119}=\) & \(\displaystyle -1 +z_{107} -z_{160}\) \\
 & \(\displaystyle -z_{161}z_{224} +z_{162} +z_{165}z_{166}\) \\[3pt]
\(r_{120}=\) & \(\displaystyle 1 +z_{151} -z_{151}z_{168}\) \\
 & \(\displaystyle -z_{168} +z_{169}\) \\[3pt]
\(r_{121}=\) & \(\displaystyle z_{167}^{2} -2z_{168} +z_{168}^{2}\) \\
 & \(\displaystyle -z_{169}^{2}\) \\[3pt]
\(r_{122}=\) & \(\displaystyle 1 -2z_{107} +z_{107}^{2}\) \\
 & \(\displaystyle -2z_{107}z_{151} +z_{170}\) \\[3pt]
\(r_{123}=\) & \(\displaystyle z_{107} -z_{107}z_{168} +z_{108}\) \\
 & \(\displaystyle -z_{167} +z_{169} +z_{170}z_{171}\) \\[3pt]
\(r_{124}=\) & \(\displaystyle 4 -4z_{107} +z_{107}^{2}\) \\
 & \(\displaystyle -2z_{107}z_{151} +2z_{151} +z_{172}\) \\[3pt]
\(r_{125}=\) & \(\displaystyle -1 +z_{107} -z_{107}z_{156}\) \\
 & \(\displaystyle +z_{109} -z_{155} +z_{156}\) \\
 & \(\displaystyle +z_{157} +z_{172}z_{173}\) \\[3pt]
\(r_{126}=\) & \(\displaystyle 1 +z_{151} -z_{151}z_{175}\) \\
 & \(\displaystyle -z_{175} +z_{176}\) \\[3pt]
\(r_{127}=\) & \(\displaystyle z_{174}^{2} -2z_{175} +z_{175}^{2}\) \\
 & \(\displaystyle -z_{176}^{2}\) \\[3pt]
\(r_{128}=\) & \(\displaystyle z_{107} -z_{107}z_{175} +z_{113}\) \\
 & \(\displaystyle +z_{170}z_{177} -z_{174} +z_{176}\) \\[3pt]
\(r_{129}=\) & \(\displaystyle -2z_{151} -z_{151}^{2} +z_{178}\) \\[3pt]
\(r_{130}=\) & \(\displaystyle -z_{107}z_{155} -z_{107}z_{157} +z_{114}\) \\
 & \(\displaystyle +z_{151}z_{157} +z_{155} -z_{156}z_{178}\) \\
 & \(\displaystyle +2z_{157} +z_{172}z_{179} +z_{178}\) \\[3pt]
\(r_{131}=\) & \(\displaystyle -z_{151}z_{181} -z_{181} +z_{182}\) \\[3pt]
\(r_{132}=\) & \(\displaystyle -1 +z_{180}^{2} +z_{181}^{2}\) \\
 & \(\displaystyle -z_{182}^{2}\) \\[3pt]
\(r_{133}=\) & \(\displaystyle -z_{107}z_{181} +z_{118} +z_{170}z_{183}\) \\
 & \(\displaystyle -z_{180} +z_{182}\) \\[3pt]
\(r_{134}=\) & \(\displaystyle -z_{151}z_{185} -z_{185} +z_{186}\) \\[3pt]
\(r_{135}=\) & \(\displaystyle -1 +z_{184}^{2} +z_{185}^{2}\) \\
 & \(\displaystyle -z_{186}^{2}\) \\[3pt]
\(r_{136}=\) & \(\displaystyle -z_{107}z_{185} +z_{121} +z_{170}z_{187}\) \\
 & \(\displaystyle -z_{184} +z_{186}\) \\[3pt]
\(r_{137}=\) & \(\displaystyle -z_{151}z_{189} -z_{189} +z_{190}\) \\[3pt]
\(r_{138}=\) & \(\displaystyle -1 +z_{188}^{2} +z_{189}^{2}\) \\
 & \(\displaystyle -z_{190}^{2}\) \\[3pt]
\(r_{139}=\) & \(\displaystyle -z_{107}z_{189} +z_{124} +z_{170}z_{191}\) \\
 & \(\displaystyle -z_{188} +z_{190}\) \\[3pt]
\(r_{140}=\) & \(\displaystyle -z_{151}z_{193} -z_{193} +z_{194}\) \\[3pt]
\(r_{141}=\) & \(\displaystyle -1 +z_{192}^{2} +z_{193}^{2}\) \\
 & \(\displaystyle -z_{194}^{2}\) \\[3pt]
\(r_{142}=\) & \(\displaystyle -z_{107}z_{193} +z_{127} +z_{170}z_{195}\) \\
 & \(\displaystyle -z_{192} +z_{194}\) \\[3pt]
\(r_{143}=\) & \(\displaystyle -1 -z_{151} -z_{151}z_{197}\) \\
 & \(\displaystyle -z_{197} +z_{198}\) \\[3pt]
\(r_{144}=\) & \(\displaystyle z_{196}^{2} +2z_{197} +z_{197}^{2}\) \\
 & \(\displaystyle -z_{198}^{2}\) \\[3pt]
\(r_{145}=\) & \(\displaystyle 1 -z_{107} -z_{107}z_{197}\) \\
 & \(\displaystyle +z_{128} +z_{172}z_{199} -z_{196}\) \\
 & \(\displaystyle +z_{197} +z_{198}\) \\[3pt]
\(r_{146}=\) & \(\displaystyle -z_{151}z_{201} -z_{201} +z_{202}\) \\[3pt]
\(r_{147}=\) & \(\displaystyle -1 +z_{200}^{2} +z_{201}^{2}\) \\
 & \(\displaystyle -z_{202}^{2}\) \\[3pt]
\(r_{148}=\) & \(\displaystyle -z_{107}z_{201} +z_{132} +z_{170}z_{203}\) \\
 & \(\displaystyle -z_{200} +z_{202}\) \\[3pt]
\(r_{149}=\) & \(\displaystyle -z_{151}z_{205} -z_{205} +z_{206}\) \\[3pt]
\(r_{150}=\) & \(\displaystyle -1 +z_{204}^{2} +z_{205}^{2}\) \\
 & \(\displaystyle -z_{206}^{2}\) \\[3pt]
\(r_{151}=\) & \(\displaystyle -z_{107}z_{205} +z_{135} +z_{170}z_{207}\) \\
 & \(\displaystyle -z_{204} +z_{206}\) \\[3pt]
\(r_{152}=\) & \(\displaystyle -z_{151}z_{209} -z_{209} +z_{210}\) \\[3pt]
\(r_{153}=\) & \(\displaystyle -1 +z_{208}^{2} +z_{209}^{2}\) \\
 & \(\displaystyle -z_{210}^{2}\) \\[3pt]
\(r_{154}=\) & \(\displaystyle -z_{107}z_{209} +z_{136} +z_{172}z_{211}\) \\
 & \(\displaystyle -z_{208} +z_{209} +z_{210}\) \\[3pt]
\(r_{155}=\) & \(\displaystyle -z_{151}z_{213} -z_{213} +z_{214}\) \\[3pt]
\(r_{156}=\) & \(\displaystyle -1 +z_{212}^{2} +z_{213}^{2}\) \\
 & \(\displaystyle -z_{214}^{2}\) \\[3pt]
\(r_{157}=\) & \(\displaystyle -z_{107}z_{213} +z_{140} +z_{170}z_{215}\) \\
 & \(\displaystyle -z_{212} +z_{214}\) \\[3pt]
\(r_{158}=\) & \(\displaystyle -z_{151}z_{217} -z_{217} +z_{218}\) \\[3pt]
\(r_{159}=\) & \(\displaystyle -1 +z_{216}^{2} +z_{217}^{2}\) \\
 & \(\displaystyle -z_{218}^{2}\) \\[3pt]
\(r_{160}=\) & \(\displaystyle -z_{107}z_{217} +z_{143} +z_{170}z_{219}\) \\
 & \(\displaystyle -z_{216} +z_{218}\) \\[3pt]
\(r_{161}=\) & \(\displaystyle -z_{151}z_{221} -z_{221} +z_{222}\) \\[3pt]
\(r_{162}=\) & \(\displaystyle -1 +z_{220}^{2} +z_{221}^{2}\) \\
 & \(\displaystyle -z_{222}^{2}\) \\[3pt]
\(r_{163}=\) & \(\displaystyle -z_{107}z_{221} +z_{147} +z_{170}z_{223}\) \\
 & \(\displaystyle -z_{220} +z_{222}\) \\[3pt]
\(r_{164}=\) & \(\displaystyle -16 +8z_{0} -14z_{0}^{2}\) \\
 & \(\displaystyle -16z_{0}z_{1} -16z_{0}z_{2} -16z_{0}z_{3}\) \\
 & \(\displaystyle -4z_{0}z_{4} -10z_{1}^{2} -16z_{1}z_{2}\) \\
 & \(\displaystyle -16z_{1}z_{3} -4z_{2} -10z_{2}^{2}\) \\
 & \(\displaystyle -16z_{2}z_{3} -8z_{3}^{2} -4z_{4}^{2}\) \\
 & \(\displaystyle -2z_{5}^{2} -2z_{6}^{2} +4z_{10}\) \\
 & \(\displaystyle -2z_{10}^{2} -2z_{11}^{2} -2z_{12}^{2}\) \\
 & \(\displaystyle -2z_{13}^{2} -2z_{14}^{2} -2z_{15}^{2}\) \\
 & \(\displaystyle -2z_{16}^{2} -2z_{17}^{2} -2z_{18}^{2}\) \\
 & \(\displaystyle -2z_{19}^{2} -2z_{105}^{2} +4z_{107}\) \\
 & \(\displaystyle -4z_{107}^{2} -2z_{108}^{2} -2z_{109}^{2}\) \\
 & \(\displaystyle -2z_{113}^{2} -2z_{114}^{2} -2z_{118}^{2}\) \\
 & \(\displaystyle -2z_{121}^{2} -2z_{124}^{2} -2z_{127}^{2}\) \\
 & \(\displaystyle -2z_{128}^{2} -2z_{132}^{2} -2z_{135}^{2}\) \\
 & \(\displaystyle -2z_{136}^{2} -2z_{140}^{2} -2z_{143}^{2}\) \\
 & \(\displaystyle -2z_{147}^{2} +z_{150} -2z_{224}^{2}\) \\[3pt]
\(r_{165}=\) & \(\displaystyle -z_{5}z_{10} -z_{5}z_{11} -z_{5}z_{12}\) \\
 & \(\displaystyle -z_{5}z_{13} -z_{5}z_{14} -z_{5}z_{15}\) \\
 & \(\displaystyle -z_{5}z_{16} -z_{5}z_{17} -z_{5}z_{18}\) \\
 & \(\displaystyle +z_{224}\) \\[3pt]
\(r_{166}=\) & \(\displaystyle -17 +z_{0} -z_{225}^{2}\) \\[3pt]
\(r_{167}=\) & \(\displaystyle z_{2} +z_{6} -z_{6}^{2}\) \\[3pt]
\(r_{168}=\) & \(\displaystyle -z_{151}z_{227} -z_{227} +z_{234}\) \\[3pt]
\(r_{169}=\) & \(\displaystyle -1 +z_{226}^{2} +z_{227}^{2}\) \\
 & \(\displaystyle -z_{234}^{2}\) \\[3pt]
\(r_{170}=\) & \(\displaystyle -z_{229} -z_{229}z_{230} +z_{235}\) \\[3pt]
\(r_{171}=\) & \(\displaystyle -1 +z_{228}^{2} +z_{229}^{2}\) \\
 & \(\displaystyle -z_{235}^{2}\) \\[3pt]
\(r_{172}=\) & \(\displaystyle -z_{221}^{2} +z_{236}\) \\[3pt]
\(r_{173}=\) & \(\displaystyle z_{227} -2z_{231}z_{236} -z_{236}\) \\[3pt]
\(r_{174}=\) & \(\displaystyle -z_{220} -z_{226}z_{232} +z_{228}\) \\[3pt]
\(r_{175}=\) & \(\displaystyle -z_{105} -4z_{221}z_{233} +z_{229}\) \\[3pt]
\(r_{176}=\) & \(\displaystyle 6z_{0}^{2} +8z_{0}z_{1} +8z_{0}z_{2}\) \\
 & \(\displaystyle +8z_{0}z_{3} +2z_{0}z_{4} -2z_{0}z_{156}\) \\
 & \(\displaystyle -4z_{0}z_{161} -2z_{0}z_{168} +5z_{1}^{2}\) \\
 & \(\displaystyle +8z_{1}z_{2} +8z_{1}z_{3} -4z_{1}z_{161}\) \\
 & \(\displaystyle -2z_{1}z_{209} +5z_{2}^{2} +8z_{2}z_{3}\) \\
 & \(\displaystyle -4z_{2}z_{161} -2z_{2}z_{197} +4z_{3}^{2}\) \\
 & \(\displaystyle -4z_{3}z_{161} +z_{4}^{2} -2z_{4}z_{168}\) \\
 & \(\displaystyle +z_{10}^{2} -2z_{10}z_{175} +z_{11}^{2}\) \\
 & \(\displaystyle -2z_{11}z_{181} +z_{12}^{2} -2z_{12}z_{185}\) \\
 & \(\displaystyle +z_{13}^{2} -2z_{13}z_{189} +z_{14}^{2}\) \\
 & \(\displaystyle -2z_{14}z_{193} +z_{15}^{2} -2z_{15}z_{201}\) \\
 & \(\displaystyle +z_{16}^{2} -2z_{16}z_{205} +z_{17}^{2}\) \\
 & \(\displaystyle -2z_{17}z_{213} +z_{18}^{2} -2z_{18}z_{217}\) \\
 & \(\displaystyle +z_{105}^{2} -2z_{105}z_{221} -z_{151}z_{237}\) \\
 & \(\displaystyle +z_{156}^{2} +z_{161}^{2} +z_{168}^{2}\) \\
 & \(\displaystyle +z_{175}^{2} +z_{181}^{2} +z_{185}^{2}\) \\
 & \(\displaystyle +z_{189}^{2} +z_{193}^{2} +z_{197}^{2}\) \\
 & \(\displaystyle +z_{201}^{2} +z_{205}^{2} +z_{209}^{2}\) \\
 & \(\displaystyle +z_{213}^{2} +z_{217}^{2} +z_{221}^{2}\) \\[3pt]
\(r_{177}=\) & \(\displaystyle -1 +z_{1} -z_{238}^{2}\) \\
 & \(\displaystyle -z_{239}^{2} -z_{240} -z_{240}^{2}\) \\[3pt]
\(r_{178}=\) & \(\displaystyle -1 -z_{3}z_{5} +z_{5}\) \\
 & \(\displaystyle +z_{6} -z_{241}^{2} -z_{242}^{2}\) \\
 & \(\displaystyle -z_{243} -z_{243}^{2}\) \\[3pt]
\(r_{179}=\) & \(\displaystyle z_{3}z_{5} -z_{6} -z_{244}^{2}\) \\
 & \(\displaystyle -z_{245}^{2} -z_{246} -z_{246}^{2}\) \\[3pt]
\(r_{180}=\) & \(\displaystyle -2 +z_{10} -z_{247}^{2}\) \\
 & \(\displaystyle -z_{248}^{2} -z_{249} -z_{249}^{2}\) \\[3pt]
\(r_{181}=\) & \(\displaystyle -1 +z_{11} -z_{250}^{2}\) \\
 & \(\displaystyle -z_{251}^{2} -z_{252} -z_{252}^{2}\) \\[3pt]
\(r_{182}=\) & \(\displaystyle -1 +z_{12} -z_{253}^{2}\) \\
 & \(\displaystyle -z_{254}^{2} -z_{255} -z_{255}^{2}\) \\[3pt]
\(r_{183}=\) & \(\displaystyle -1 +z_{13} -z_{256}^{2}\) \\
 & \(\displaystyle -z_{257}^{2} -z_{258} -z_{258}^{2}\) \\[3pt]
\(r_{184}=\) & \(\displaystyle -1 +z_{14} -z_{259}^{2}\) \\
 & \(\displaystyle -z_{260}^{2} -z_{261} -z_{261}^{2}\) \\[3pt]
\(r_{185}=\) & \(\displaystyle -1 +z_{15} -z_{262}^{2}\) \\
 & \(\displaystyle -z_{263}^{2} -z_{264} -z_{264}^{2}\) \\[3pt]
\(r_{186}=\) & \(\displaystyle -1 +z_{16} -z_{265}^{2}\) \\
 & \(\displaystyle -z_{266}^{2} -z_{267} -z_{267}^{2}\) \\[3pt]
\(r_{187}=\) & \(\displaystyle -1 +z_{17} -z_{268}^{2}\) \\
 & \(\displaystyle -z_{269}^{2} -z_{270} -z_{270}^{2}\) \\[3pt]
\(r_{188}=\) & \(\displaystyle -1 +z_{18} -z_{271}^{2}\) \\
 & \(\displaystyle -z_{272}^{2} -z_{273} -z_{273}^{2}\) \\[3pt]
\(r_{189}=\) & \(\displaystyle -2 -z_{5} +z_{107}\) \\
 & \(\displaystyle -z_{274}^{2} -z_{275}^{2} -z_{276}\) \\
 & \(\displaystyle -z_{276}^{2}\) \\[3pt]
\(r_{190}=\) & \(\displaystyle z_{109}z_{111} -z_{111} -z_{111}^{2}\) \\
 & \(\displaystyle -z_{277}^{2} -z_{278}^{2} -z_{279}\) \\
 & \(\displaystyle -z_{279}^{2}\) \\[3pt]
\(r_{191}=\) & \(\displaystyle z_{114}z_{116} -z_{116} -z_{116}^{2}\) \\
 & \(\displaystyle -z_{280}^{2} -z_{281}^{2} -z_{282}\) \\
 & \(\displaystyle -z_{282}^{2}\) \\[3pt]
\(r_{192}=\) & \(\displaystyle z_{114}z_{120} -z_{120} -z_{120}^{2}\) \\
 & \(\displaystyle -z_{283}^{2} -z_{284}^{2} -z_{285}\) \\
 & \(\displaystyle -z_{285}^{2}\) \\[3pt]
\(r_{193}=\) & \(\displaystyle z_{114}z_{123} -z_{123} -z_{123}^{2}\) \\
 & \(\displaystyle -z_{286}^{2} -z_{287}^{2} -z_{288}\) \\
 & \(\displaystyle -z_{288}^{2}\) \\[3pt]
\(r_{194}=\) & \(\displaystyle z_{114}z_{126} -z_{126} -z_{126}^{2}\) \\
 & \(\displaystyle -z_{289}^{2} -z_{290}^{2} -z_{291}\) \\
 & \(\displaystyle -z_{291}^{2}\) \\[3pt]
\(r_{195}=\) & \(\displaystyle z_{128}z_{130} -z_{130} -z_{130}^{2}\) \\
 & \(\displaystyle -z_{292}^{2} -z_{293}^{2} -z_{294}\) \\
 & \(\displaystyle -z_{294}^{2}\) \\[3pt]
\(r_{196}=\) & \(\displaystyle z_{128}z_{134} -z_{134} -z_{134}^{2}\) \\
 & \(\displaystyle -z_{295}^{2} -z_{296}^{2} -z_{297}\) \\
 & \(\displaystyle -z_{297}^{2}\) \\[3pt]
\(r_{197}=\) & \(\displaystyle z_{136}z_{138} -z_{138} -z_{138}^{2}\) \\
 & \(\displaystyle -z_{298}^{2} -z_{299}^{2} -z_{300}\) \\
 & \(\displaystyle -z_{300}^{2}\) \\[3pt]
\(r_{198}=\) & \(\displaystyle z_{136}z_{142} -z_{142} -z_{142}^{2}\) \\
 & \(\displaystyle -z_{301}^{2} -z_{302}^{2} -z_{303}\) \\
 & \(\displaystyle -z_{303}^{2}\) \\[3pt]
\(r_{199}=\) & \(\displaystyle z_{109}z_{145} -z_{145} -z_{145}^{2}\) \\
 & \(\displaystyle -z_{304}^{2} -z_{305}^{2} -z_{306}\) \\
 & \(\displaystyle -z_{306}^{2}\) \\[3pt]
\(r_{200}=\) & \(\displaystyle z_{114}z_{149} -z_{149} -z_{149}^{2}\) \\
 & \(\displaystyle -z_{307}^{2} -z_{308}^{2} -z_{309}\) \\
 & \(\displaystyle -z_{309}^{2}\) \\[3pt]
\(r_{201}=\) & \(\displaystyle -5 -z_{150}^{2} +2z_{156}\) \\
 & \(\displaystyle -z_{156}^{2} -z_{161}^{2} +2z_{168}\) \\
 & \(\displaystyle -z_{168}^{2} +2z_{175} -z_{175}^{2}\) \\
 & \(\displaystyle -z_{181}^{2} -z_{185}^{2} -z_{189}^{2}\) \\
 & \(\displaystyle -z_{193}^{2} -2z_{197} -z_{197}^{2}\) \\
 & \(\displaystyle -z_{201}^{2} -z_{205}^{2} -z_{209}^{2}\) \\
 & \(\displaystyle -z_{213}^{2} -z_{217}^{2} +z_{221}^{2}\) \\
 & \(\displaystyle -z_{310}^{2} -z_{311}^{2} -z_{312}\) \\
 & \(\displaystyle -z_{312}^{2}\) \\[3pt]
\(r_{202}=\) & \(\displaystyle 1 -z_{226}^{2} +z_{313}\) \\[3pt]
\(r_{203}=\) & \(\displaystyle z_{151}z_{313} +z_{230}\) \\[3pt]
\end{longtable}
\endgroup


\section{The complete natural-number residual list}
\label{app:natural}
All 233 variables \(x_0,\ldots,x_{232}\) range independently over
\(\N\). The equation is \(\sum_{j=0}^{181}s_j^2=0\), using every
row below. The variable-name map in the artifact is explanatory only;
the indexed list is already a complete polynomial specification.
% Generated from natural/equation.json; all residuals are printed.
% All x_i independently range over the nonnegative integers.
\begingroup
\small
\setlength{\tabcolsep}{2pt}
\setlength{\LTleft}{0pt}
\begin{longtable}{@{}r@{\;}l@{}}
\(s_{0}=\) & \(\displaystyle -3 -4x_{2} -x_{2}^{2}\) \\
 & \(\displaystyle +x_{3}\) \\[3pt]
\(s_{1}=\) & \(\displaystyle 1 +x_{4} -x_{5}\) \\
 & \(\displaystyle -x_{6}\) \\[3pt]
\(s_{2}=\) & \(\displaystyle -39 -3x_{0} -x_{0}x_{2}\) \\
 & \(\displaystyle -x_{0}x_{8} -13x_{2} -13x_{8}\) \\
 & \(\displaystyle +x_{10}\) \\[3pt]
\(s_{3}=\) & \(\displaystyle -3 +x_{1} -x_{2}\) \\
 & \(\displaystyle -x_{2}x_{7} -4x_{7} -x_{7}x_{8}\) \\[3pt]
\(s_{4}=\) & \(\displaystyle -1 +x_{1} -x_{9}\) \\
 & \(\displaystyle -x_{9}x_{10}\) \\[3pt]
\(s_{5}=\) & \(\displaystyle 1 -x_{3} -x_{4}x_{21}\) \\
 & \(\displaystyle +x_{19} -2x_{21}\) \\[3pt]
\(s_{6}=\) & \(\displaystyle -6 -5x_{11} -x_{11}^{2}\) \\
 & \(\displaystyle -2x_{11}x_{12} -5x_{12} -x_{12}^{2}\) \\
 & \(\displaystyle +x_{22}\) \\[3pt]
\(s_{7}=\) & \(\displaystyle -72 -68x_{12} -16x_{12}^{2}\) \\
 & \(\displaystyle -16x_{12}x_{22} -34x_{22} -4x_{22}^{2}\) \\
 & \(\displaystyle +x_{23}\) \\[3pt]
\(s_{8}=\) & \(\displaystyle -30 -11x_{11} -x_{11}^{2}\) \\
 & \(\displaystyle -4x_{11}x_{23} -22x_{23} -4x_{23}^{2}\) \\
 & \(\displaystyle +x_{24}\) \\[3pt]
\(s_{9}=\) & \(\displaystyle -380 -312x_{23} -64x_{23}^{2}\) \\
 & \(\displaystyle -32x_{23}x_{24} -78x_{24} -4x_{24}^{2}\) \\
 & \(\displaystyle +x_{25}\) \\[3pt]
\(s_{10}=\) & \(\displaystyle -6 -5x_{12} -x_{12}^{2}\) \\
 & \(\displaystyle -2x_{12}x_{13} -5x_{13} -x_{13}^{2}\) \\
 & \(\displaystyle +x_{26}\) \\[3pt]
\(s_{11}=\) & \(\displaystyle -6 -5x_{11} -x_{11}^{2}\) \\
 & \(\displaystyle -2x_{11}x_{13} -5x_{13} -x_{13}^{2}\) \\
 & \(\displaystyle +x_{27}\) \\[3pt]
\(s_{12}=\) & \(\displaystyle -414 -22x_{11} -2x_{11}^{2}\) \\
 & \(\displaystyle -16x_{11}x_{13} -8x_{11}x_{26} -200x_{12}\) \\
 & \(\displaystyle -32x_{12}^{2} -64x_{12}x_{13} -32x_{12}x_{22}\) \\
 & \(\displaystyle -32x_{12}x_{27} -320x_{13} -64x_{13}^{2}\) \\
 & \(\displaystyle -32x_{13}x_{22} -32x_{13}x_{26} -32x_{13}x_{27}\) \\
 & \(\displaystyle -100x_{22} -8x_{22}^{2} -16x_{22}x_{27}\) \\
 & \(\displaystyle -52x_{26} -8x_{26}^{2} -108x_{27}\) \\
 & \(\displaystyle -8x_{27}^{2} +x_{28}\) \\[3pt]
\(s_{13}=\) & \(\displaystyle -1344 -800x_{12} -128x_{12}^{2}\) \\
 & \(\displaystyle -256x_{12}x_{13} -128x_{12}x_{22} -128x_{12}x_{27}\) \\
 & \(\displaystyle -864x_{13} -128x_{13}^{2} -128x_{13}x_{22}\) \\
 & \(\displaystyle -128x_{13}x_{27} -400x_{22} -32x_{22}^{2}\) \\
 & \(\displaystyle -64x_{22}x_{27} -432x_{27} -32x_{27}^{2}\) \\
 & \(\displaystyle -2x_{28} -2x_{28}^{2} +x_{29}\) \\[3pt]
\(s_{14}=\) & \(\displaystyle -247 -168x_{11} -32x_{11}^{2}\) \\
 & \(\displaystyle -64x_{11}x_{12} -32x_{11}x_{26} -32x_{11}x_{27}\) \\
 & \(\displaystyle -188x_{12} -32x_{12}^{2} -32x_{12}x_{26}\) \\
 & \(\displaystyle -32x_{12}x_{27} -2x_{22} -92x_{26}\) \\
 & \(\displaystyle -8x_{26}^{2} -16x_{26}x_{27} -84x_{27}\) \\
 & \(\displaystyle -8x_{27}^{2} +x_{30}\) \\[3pt]
\(s_{15}=\) & \(\displaystyle -960 -672x_{11} -128x_{11}^{2}\) \\
 & \(\displaystyle -256x_{11}x_{12} -128x_{11}x_{26} -128x_{11}x_{27}\) \\
 & \(\displaystyle -736x_{12} -128x_{12}^{2} -128x_{12}x_{26}\) \\
 & \(\displaystyle -128x_{12}x_{27} -368x_{26} -32x_{26}^{2}\) \\
 & \(\displaystyle -64x_{26}x_{27} -336x_{27} -32x_{27}^{2}\) \\
 & \(\displaystyle -2x_{30} -2x_{30}^{2} +x_{31}\) \\[3pt]
\(s_{16}=\) & \(\displaystyle -276 -30x_{11} -2x_{11}^{2}\) \\
 & \(\displaystyle -16x_{11}x_{13} -8x_{11}x_{26} -30x_{12}\) \\
 & \(\displaystyle -2x_{12}^{2} -16x_{12}x_{13} -8x_{12}x_{27}\) \\
 & \(\displaystyle -272x_{13} -64x_{13}^{2} -32x_{13}x_{26}\) \\
 & \(\displaystyle -32x_{13}x_{27} -68x_{26} -8x_{26}^{2}\) \\
 & \(\displaystyle -68x_{27} -8x_{27}^{2} +x_{32}\) \\[3pt]
\(s_{17}=\) & \(\displaystyle -552 -120x_{12} -8x_{12}^{2}\) \\
 & \(\displaystyle -64x_{12}x_{13} -32x_{12}x_{27} -544x_{13}\) \\
 & \(\displaystyle -128x_{13}^{2} -128x_{13}x_{27} -272x_{27}\) \\
 & \(\displaystyle -32x_{27}^{2} -2x_{32} -2x_{32}^{2}\) \\
 & \(\displaystyle +x_{33}\) \\[3pt]
\(s_{18}=\) & \(\displaystyle -247 -4x_{11} -352x_{12}\) \\
 & \(\displaystyle -128x_{12}^{2} -64x_{12}x_{22} -64x_{12}x_{26}\) \\
 & \(\displaystyle -92x_{22} -8x_{22}^{2} -16x_{22}x_{26}\) \\
 & \(\displaystyle -84x_{26} -8x_{26}^{2} -2x_{27}\) \\
 & \(\displaystyle +x_{34}\) \\[3pt]
\(s_{19}=\) & \(\displaystyle -960 -1408x_{12} -512x_{12}^{2}\) \\
 & \(\displaystyle -256x_{12}x_{22} -256x_{12}x_{26} -368x_{22}\) \\
 & \(\displaystyle -32x_{22}^{2} -64x_{22}x_{26} -336x_{26}\) \\
 & \(\displaystyle -32x_{26}^{2} -2x_{34} -2x_{34}^{2}\) \\
 & \(\displaystyle +x_{35}\) \\[3pt]
\(s_{20}=\) & \(\displaystyle -56 -60x_{12} -16x_{12}^{2}\) \\
 & \(\displaystyle -16x_{12}x_{22} -30x_{22} -4x_{22}^{2}\) \\
 & \(\displaystyle +x_{36}\) \\[3pt]
\(s_{21}=\) & \(\displaystyle -30 -11x_{11} -x_{11}^{2}\) \\
 & \(\displaystyle -4x_{11}x_{36} -22x_{36} -4x_{36}^{2}\) \\
 & \(\displaystyle +x_{37}\) \\[3pt]
\(s_{22}=\) & \(\displaystyle -380 -312x_{36} -64x_{36}^{2}\) \\
 & \(\displaystyle -32x_{36}x_{37} -78x_{37} -4x_{37}^{2}\) \\
 & \(\displaystyle +x_{38}\) \\[3pt]
\(s_{23}=\) & \(\displaystyle -30 -11x_{12} -x_{12}^{2}\) \\
 & \(\displaystyle -4x_{12}x_{38} -22x_{38} -4x_{38}^{2}\) \\
 & \(\displaystyle +x_{39}\) \\[3pt]
\(s_{24}=\) & \(\displaystyle -349 -420x_{12} -128x_{12}^{2}\) \\
 & \(\displaystyle -128x_{12}x_{22} -208x_{22} -32x_{22}^{2}\) \\
 & \(\displaystyle -2x_{26} +x_{40}\) \\[3pt]
\(s_{25}=\) & \(\displaystyle -1368 -1664x_{12} -512x_{12}^{2}\) \\
 & \(\displaystyle -512x_{12}x_{22} -832x_{22} -128x_{22}^{2}\) \\
 & \(\displaystyle -2x_{40} -2x_{40}^{2} +x_{41}\) \\[3pt]
\(s_{26}=\) & \(\displaystyle -56 -60x_{11} -16x_{11}^{2}\) \\
 & \(\displaystyle -16x_{11}x_{27} -8x_{11}x_{41} -30x_{27}\) \\
 & \(\displaystyle -4x_{27}^{2} -4x_{27}x_{41} -15x_{41}\) \\
 & \(\displaystyle -x_{41}^{2} +x_{42}\) \\[3pt]
\(s_{27}=\) & \(\displaystyle -6 -5x_{13} -x_{13}^{2}\) \\
 & \(\displaystyle -2x_{13}x_{14} -5x_{14} -x_{14}^{2}\) \\
 & \(\displaystyle +x_{43}\) \\[3pt]
\(s_{28}=\) & \(\displaystyle -224 -26x_{11} -2x_{11}^{2}\) \\
 & \(\displaystyle -16x_{11}x_{14} -8x_{11}x_{43} -288x_{13}\) \\
 & \(\displaystyle -128x_{13}^{2} -64x_{13}x_{26} -64x_{13}x_{27}\) \\
 & \(\displaystyle -88x_{14} -32x_{14}^{2} -32x_{14}x_{43}\) \\
 & \(\displaystyle -76x_{26} -8x_{26}^{2} -16x_{26}x_{27}\) \\
 & \(\displaystyle -68x_{27} -8x_{27}^{2} -44x_{43}\) \\
 & \(\displaystyle -8x_{43}^{2} +x_{44}\) \\[3pt]
\(s_{29}=\) & \(\displaystyle -640 -1152x_{13} -512x_{13}^{2}\) \\
 & \(\displaystyle -256x_{13}x_{26} -256x_{13}x_{27} -304x_{26}\) \\
 & \(\displaystyle -32x_{26}^{2} -64x_{26}x_{27} -272x_{27}\) \\
 & \(\displaystyle -32x_{27}^{2} -2x_{44} -2x_{44}^{2}\) \\
 & \(\displaystyle +x_{45}\) \\[3pt]
\(s_{30}=\) & \(\displaystyle -241 -168x_{11} -32x_{11}^{2}\) \\
 & \(\displaystyle -64x_{11}x_{12} -32x_{11}x_{26} -32x_{11}x_{27}\) \\
 & \(\displaystyle -184x_{12} -32x_{12}^{2} -32x_{12}x_{26}\) \\
 & \(\displaystyle -32x_{12}x_{27} -x_{13} -92x_{26}\) \\
 & \(\displaystyle -8x_{26}^{2} -16x_{26}x_{27} -84x_{27}\) \\
 & \(\displaystyle -8x_{27}^{2} +x_{46}\) \\[3pt]
\(s_{31}=\) & \(\displaystyle -42 -52x_{13} -16x_{13}^{2}\) \\
 & \(\displaystyle -16x_{13}x_{26} -26x_{26} -4x_{26}^{2}\) \\
 & \(\displaystyle +x_{47}\) \\[3pt]
\(s_{32}=\) & \(\displaystyle -30 -11x_{13} -x_{13}^{2}\) \\
 & \(\displaystyle -4x_{13}x_{47} -22x_{47} -4x_{47}^{2}\) \\
 & \(\displaystyle +x_{48}\) \\[3pt]
\(s_{33}=\) & \(\displaystyle -306 -35x_{12} -x_{12}^{2}\) \\
 & \(\displaystyle -16x_{12}x_{47} -4x_{12}x_{48} -280x_{47}\) \\
 & \(\displaystyle -64x_{47}^{2} -32x_{47}x_{48} -70x_{48}\) \\
 & \(\displaystyle -4x_{48}^{2} +x_{49}\) \\[3pt]
\(s_{34}=\) & \(\displaystyle -4556 -4320x_{47} -1024x_{47}^{2}\) \\
 & \(\displaystyle -512x_{47}x_{48} -128x_{47}x_{49} -1080x_{48}\) \\
 & \(\displaystyle -64x_{48}^{2} -32x_{48}x_{49} -270x_{49}\) \\
 & \(\displaystyle -4x_{49}^{2} +x_{50}\) \\[3pt]
\(s_{35}=\) & \(\displaystyle -30 -11x_{13} -x_{13}^{2}\) \\
 & \(\displaystyle -4x_{13}x_{50} -22x_{50} -4x_{50}^{2}\) \\
 & \(\displaystyle +x_{51}\) \\[3pt]
\(s_{36}=\) & \(\displaystyle -306 -35x_{12} -x_{12}^{2}\) \\
 & \(\displaystyle -16x_{12}x_{50} -4x_{12}x_{51} -280x_{50}\) \\
 & \(\displaystyle -64x_{50}^{2} -32x_{50}x_{51} -70x_{51}\) \\
 & \(\displaystyle -4x_{51}^{2} +x_{52}\) \\[3pt]
\(s_{37}=\) & \(\displaystyle -1028 -672x_{11} -128x_{11}^{2}\) \\
 & \(\displaystyle -256x_{11}x_{12} -128x_{11}x_{26} -128x_{11}x_{27}\) \\
 & \(\displaystyle -736x_{12} -128x_{12}^{2} -128x_{12}x_{26}\) \\
 & \(\displaystyle -128x_{12}x_{27} -368x_{26} -32x_{26}^{2}\) \\
 & \(\displaystyle -64x_{26}x_{27} -336x_{27} -32x_{27}^{2}\) \\
 & \(\displaystyle -2x_{46} -2x_{46}^{2} -32x_{50}\) \\
 & \(\displaystyle -8x_{51} -2x_{52} +x_{53}\) \\[3pt]
\(s_{38}=\) & \(\displaystyle -264 -128x_{50} -32x_{51}\) \\
 & \(\displaystyle -8x_{52} -2x_{53} -2x_{53}^{2}\) \\
 & \(\displaystyle +x_{54}\) \\[3pt]
\(s_{39}=\) & \(\displaystyle -2 -3x_{11} -x_{11}^{2}\) \\
 & \(\displaystyle -2x_{11}x_{54} -3x_{54} -x_{54}^{2}\) \\
 & \(\displaystyle +x_{55}\) \\[3pt]
\(s_{40}=\) & \(\displaystyle -12 -28x_{54} -16x_{54}^{2}\) \\
 & \(\displaystyle -16x_{54}x_{55} -14x_{55} -4x_{55}^{2}\) \\
 & \(\displaystyle +x_{56}\) \\[3pt]
\(s_{41}=\) & \(\displaystyle -30 -11x_{13} -x_{13}^{2}\) \\
 & \(\displaystyle -4x_{13}x_{56} -22x_{56} -4x_{56}^{2}\) \\
 & \(\displaystyle +x_{57}\) \\[3pt]
\(s_{42}=\) & \(\displaystyle -306 -35x_{11} -x_{11}^{2}\) \\
 & \(\displaystyle -16x_{11}x_{56} -4x_{11}x_{57} -280x_{56}\) \\
 & \(\displaystyle -64x_{56}^{2} -32x_{56}x_{57} -70x_{57}\) \\
 & \(\displaystyle -4x_{57}^{2} +x_{58}\) \\[3pt]
\(s_{43}=\) & \(\displaystyle -4556 -4320x_{56} -1024x_{56}^{2}\) \\
 & \(\displaystyle -512x_{56}x_{57} -128x_{56}x_{58} -1080x_{57}\) \\
 & \(\displaystyle -64x_{57}^{2} -32x_{57}x_{58} -270x_{58}\) \\
 & \(\displaystyle -4x_{58}^{2} +x_{59}\) \\[3pt]
\(s_{44}=\) & \(\displaystyle -90 -76x_{13} -16x_{13}^{2}\) \\
 & \(\displaystyle -16x_{13}x_{26} -16x_{13}x_{59} -38x_{26}\) \\
 & \(\displaystyle -4x_{26}^{2} -8x_{26}x_{59} -38x_{59}\) \\
 & \(\displaystyle -4x_{59}^{2} +x_{60}\) \\[3pt]
\(s_{45}=\) & \(\displaystyle -4556 -4320x_{50} -1024x_{50}^{2}\) \\
 & \(\displaystyle -512x_{50}x_{51} -128x_{50}x_{52} -1080x_{51}\) \\
 & \(\displaystyle -64x_{51}^{2} -32x_{51}x_{52} -270x_{52}\) \\
 & \(\displaystyle -4x_{52}^{2} +x_{61}\) \\[3pt]
\(s_{46}=\) & \(\displaystyle -682 -70x_{12} -2x_{12}^{2}\) \\
 & \(\displaystyle -32x_{12}x_{59} -8x_{12}x_{60} -592x_{59}\) \\
 & \(\displaystyle -128x_{59}^{2} -64x_{59}x_{60} -148x_{60}\) \\
 & \(\displaystyle -8x_{60}^{2} -2x_{61} +x_{62}\) \\[3pt]
\(s_{47}=\) & \(\displaystyle -x_{62} -x_{62}^{2} +x_{63}\) \\[3pt]
\(s_{48}=\) & \(\displaystyle -42 -52x_{12} -16x_{12}^{2}\) \\
 & \(\displaystyle -8x_{12}x_{13} -16x_{12}x_{22} -13x_{13}\) \\
 & \(\displaystyle -x_{13}^{2} -4x_{13}x_{22} -26x_{22}\) \\
 & \(\displaystyle -4x_{22}^{2} +x_{64}\) \\[3pt]
\(s_{49}=\) & \(\displaystyle -2278 -2264x_{12} -800x_{12}^{2}\) \\
 & \(\displaystyle -800x_{12}x_{22} -160x_{12}x_{64} -70x_{13}\) \\
 & \(\displaystyle -2x_{13}^{2} -32x_{13}x_{61} -8x_{13}x_{63}\) \\
 & \(\displaystyle -1132x_{22} -200x_{22}^{2} -80x_{22}x_{64}\) \\
 & \(\displaystyle -592x_{61} -128x_{61}^{2} -64x_{61}x_{63}\) \\
 & \(\displaystyle -148x_{63} -8x_{63}^{2} -228x_{64}\) \\
 & \(\displaystyle -8x_{64}^{2} +x_{65}\) \\[3pt]
\(s_{50}=\) & \(\displaystyle -2712 -280x_{13} -8x_{13}^{2}\) \\
 & \(\displaystyle -128x_{13}x_{61} -32x_{13}x_{63} -2368x_{61}\) \\
 & \(\displaystyle -512x_{61}^{2} -256x_{61}x_{63} -592x_{63}\) \\
 & \(\displaystyle -32x_{63}^{2} -2x_{65} -2x_{65}^{2}\) \\
 & \(\displaystyle +x_{66}\) \\[3pt]
\(s_{51}=\) & \(\displaystyle -2 -3x_{66} -x_{66}^{2}\) \\
 & \(\displaystyle +x_{67}\) \\[3pt]
\(s_{52}=\) & \(\displaystyle -2 -x_{2} -x_{2}x_{16}\) \\
 & \(\displaystyle -2x_{16} +x_{68}\) \\[3pt]
\(s_{53}=\) & \(\displaystyle -x_{15}^{2} +2x_{15}x_{16} -x_{16}^{2}\) \\
 & \(\displaystyle +x_{69}\) \\[3pt]
\(s_{54}=\) & \(\displaystyle -2x_{25}x_{68} -4x_{68} +x_{70}\) \\[3pt]
\(s_{55}=\) & \(\displaystyle -x_{15}^{2} +2x_{15}x_{16} +2x_{15}x_{70}\) \\
 & \(\displaystyle -x_{16}^{2} -2x_{16}x_{70} -x_{70}^{2}\) \\
 & \(\displaystyle +x_{71}\) \\[3pt]
\(s_{56}=\) & \(\displaystyle -x_{29}x_{68} +x_{72}\) \\[3pt]
\(s_{57}=\) & \(\displaystyle -x_{15}^{2} +2x_{15}x_{16} +2x_{15}x_{72}\) \\
 & \(\displaystyle -x_{16}^{2} -2x_{16}x_{72} -x_{72}^{2}\) \\
 & \(\displaystyle +x_{73}\) \\[3pt]
\(s_{58}=\) & \(\displaystyle -x_{31}x_{68} +x_{74}\) \\[3pt]
\(s_{59}=\) & \(\displaystyle -x_{15}^{2} +2x_{15}x_{16} +2x_{15}x_{74}\) \\
 & \(\displaystyle -x_{16}^{2} -2x_{16}x_{74} -x_{74}^{2}\) \\
 & \(\displaystyle +x_{75}\) \\[3pt]
\(s_{60}=\) & \(\displaystyle -x_{33}x_{68} +x_{76}\) \\[3pt]
\(s_{61}=\) & \(\displaystyle -x_{15}^{2} +2x_{15}x_{16} +2x_{15}x_{76}\) \\
 & \(\displaystyle -x_{16}^{2} -2x_{16}x_{76} -x_{76}^{2}\) \\
 & \(\displaystyle +x_{77}\) \\[3pt]
\(s_{62}=\) & \(\displaystyle -x_{35}x_{68} +x_{78}\) \\[3pt]
\(s_{63}=\) & \(\displaystyle -x_{15}^{2} +2x_{15}x_{16} +2x_{15}x_{78}\) \\
 & \(\displaystyle -x_{16}^{2} -2x_{16}x_{78} -x_{78}^{2}\) \\
 & \(\displaystyle +x_{79}\) \\[3pt]
\(s_{64}=\) & \(\displaystyle -4x_{12}x_{68} -2x_{39}x_{68} -4x_{68}\) \\
 & \(\displaystyle +x_{80}\) \\[3pt]
\(s_{65}=\) & \(\displaystyle -x_{15}^{2} +2x_{15}x_{16} +2x_{15}x_{80}\) \\
 & \(\displaystyle -x_{16}^{2} -2x_{16}x_{80} -x_{80}^{2}\) \\
 & \(\displaystyle +x_{81}\) \\[3pt]
\(s_{66}=\) & \(\displaystyle -4x_{41}x_{68} -2x_{42}x_{68} +x_{82}\) \\[3pt]
\(s_{67}=\) & \(\displaystyle -x_{15}^{2} +2x_{15}x_{16} +2x_{15}x_{82}\) \\
 & \(\displaystyle -x_{16}^{2} -2x_{16}x_{82} -x_{82}^{2}\) \\
 & \(\displaystyle +x_{83}\) \\[3pt]
\(s_{68}=\) & \(\displaystyle -x_{45}x_{68} +x_{84}\) \\[3pt]
\(s_{69}=\) & \(\displaystyle -x_{15}^{2} +2x_{15}x_{16} +2x_{15}x_{84}\) \\
 & \(\displaystyle -x_{16}^{2} -2x_{16}x_{84} -x_{84}^{2}\) \\
 & \(\displaystyle +x_{85}\) \\[3pt]
\(s_{70}=\) & \(\displaystyle -x_{11}^{2} +2x_{11}x_{13} -x_{13}^{2}\) \\
 & \(\displaystyle +x_{86}\) \\[3pt]
\(s_{71}=\) & \(\displaystyle -x_{12}^{2} +2x_{12}x_{13} -x_{13}^{2}\) \\
 & \(\displaystyle +x_{87}\) \\[3pt]
\(s_{72}=\) & \(\displaystyle -x_{86}x_{87} +x_{88}\) \\[3pt]
\(s_{73}=\) & \(\displaystyle -4 -4x_{14} -x_{14}^{2}\) \\
 & \(\displaystyle +x_{89}\) \\[3pt]
\(s_{74}=\) & \(\displaystyle -x_{88}^{2} +2x_{88}x_{89} -x_{89}^{2}\) \\
 & \(\displaystyle +x_{90}\) \\[3pt]
\(s_{75}=\) & \(\displaystyle -2x_{67}x_{68} -4x_{68} +x_{91}\) \\[3pt]
\(s_{76}=\) & \(\displaystyle -x_{15}^{2} +2x_{15}x_{16} +2x_{15}x_{91}\) \\
 & \(\displaystyle -x_{16}^{2} -2x_{16}x_{91} -x_{91}^{2}\) \\
 & \(\displaystyle +x_{92}\) \\[3pt]
\(s_{77}=\) & \(\displaystyle -2 -x_{2} -x_{2}x_{12}\) \\
 & \(\displaystyle -2x_{12} +x_{93}\) \\[3pt]
\(s_{78}=\) & \(\displaystyle -1 -x_{16} -x_{16}x_{93}\) \\
 & \(\displaystyle -x_{93} +x_{94}\) \\[3pt]
\(s_{79}=\) & \(\displaystyle -2 -x_{2} -x_{2}x_{11}\) \\
 & \(\displaystyle -2x_{11} +x_{95}\) \\[3pt]
\(s_{80}=\) & \(\displaystyle -8 -4x_{2} -4x_{2}x_{12}\) \\
 & \(\displaystyle -2x_{2}x_{22} -8x_{12} -4x_{22}\) \\
 & \(\displaystyle +x_{96}\) \\[3pt]
\(s_{81}=\) & \(\displaystyle -1 -x_{15} -x_{15}x_{95}\) \\
 & \(\displaystyle -x_{95} +x_{97}\) \\[3pt]
\(s_{82}=\) & \(\displaystyle -x_{96}x_{97} -x_{97} +x_{98}\) \\[3pt]
\(s_{83}=\) & \(\displaystyle -x_{94}^{2} +2x_{94}x_{98} -x_{98}^{2}\) \\
 & \(\displaystyle +x_{99}\) \\[3pt]
\(s_{84}=\) & \(\displaystyle -12 -6x_{2} -4x_{2}x_{12}\) \\
 & \(\displaystyle -2x_{2}x_{22} -8x_{12} -4x_{22}\) \\
 & \(\displaystyle +x_{100}\) \\[3pt]
\(s_{85}=\) & \(\displaystyle -1 -x_{16} -x_{16}x_{100}\) \\
 & \(\displaystyle -x_{100} +x_{101}\) \\[3pt]
\(s_{86}=\) & \(\displaystyle -1 -x_{15} -x_{15}x_{93}\) \\
 & \(\displaystyle -x_{93} +x_{102}\) \\[3pt]
\(s_{87}=\) & \(\displaystyle -x_{101}^{2} +2x_{101}x_{102} -x_{102}^{2}\) \\
 & \(\displaystyle +x_{103}\) \\[3pt]
\(s_{88}=\) & \(\displaystyle -x_{69}x_{71} +x_{104}\) \\[3pt]
\(s_{89}=\) & \(\displaystyle -x_{73}x_{104} +x_{105}\) \\[3pt]
\(s_{90}=\) & \(\displaystyle -x_{75}x_{105} +x_{106}\) \\[3pt]
\(s_{91}=\) & \(\displaystyle -x_{77}x_{106} +x_{107}\) \\[3pt]
\(s_{92}=\) & \(\displaystyle -x_{79}x_{107} +x_{108}\) \\[3pt]
\(s_{93}=\) & \(\displaystyle -x_{81}x_{108} +x_{109}\) \\[3pt]
\(s_{94}=\) & \(\displaystyle -x_{83}x_{109} +x_{110}\) \\[3pt]
\(s_{95}=\) & \(\displaystyle -x_{85}x_{110} +x_{111}\) \\[3pt]
\(s_{96}=\) & \(\displaystyle -x_{90}x_{111} -x_{92}x_{111} +x_{112}\) \\[3pt]
\(s_{97}=\) & \(\displaystyle -x_{99}x_{112} +x_{113}\) \\[3pt]
\(s_{98}=\) & \(\displaystyle -6 -2x_{2} -x_{2}x_{19}\) \\
 & \(\displaystyle -2x_{8} -x_{8}x_{19} -3x_{19}\) \\
 & \(\displaystyle +x_{114}\) \\[3pt]
\(s_{99}=\) & \(\displaystyle -1 -x_{17} -x_{17}x_{114}\) \\
 & \(\displaystyle -x_{114} +x_{115}\) \\[3pt]
\(s_{100}=\) & \(\displaystyle -4 -x_{2} -x_{2}x_{18}\) \\
 & \(\displaystyle -x_{8} -x_{8}x_{18} -4x_{18}\) \\
 & \(\displaystyle -x_{18}x_{114} -x_{114} +x_{116}\) \\[3pt]
\(s_{101}=\) & \(\displaystyle 2 -4x_{1} +2x_{1}^{2}\) \\
 & \(\displaystyle -2x_{1}x_{15} -2x_{1}x_{16} -2x_{1}x_{115}\) \\
 & \(\displaystyle -2x_{1}x_{116} +2x_{15} +x_{15}^{2}\) \\
 & \(\displaystyle +2x_{15}x_{115} +2x_{16} +x_{16}^{2}\) \\
 & \(\displaystyle +2x_{16}x_{116} -x_{20}x_{117} +x_{103}x_{113}\) \\
 & \(\displaystyle +2x_{115} +x_{115}^{2} +2x_{116}\) \\
 & \(\displaystyle +x_{116}^{2}\) \\[3pt]
\(s_{102}=\) & \(\displaystyle -9 -x_{11} -x_{12}\) \\
 & \(\displaystyle -x_{13} -x_{14} -x_{15}\) \\
 & \(\displaystyle -x_{16} -x_{17} -x_{18}\) \\
 & \(\displaystyle -x_{19} +10x_{20}^{2} -x_{118}\) \\[3pt]
\(s_{103}=\) & \(\displaystyle -2 -x_{4} -x_{4}x_{169}\) \\
 & \(\displaystyle -x_{4}x_{170} -3x_{169} -x_{169}x_{171}\) \\
 & \(\displaystyle -3x_{170} -x_{170}x_{171} -x_{171}\) \\
 & \(\displaystyle +x_{176} +x_{177}\) \\[3pt]
\(s_{104}=\) & \(\displaystyle -x_{119} -x_{119}x_{122} +x_{123}\) \\[3pt]
\(s_{105}=\) & \(\displaystyle -x_{119}x_{123} -x_{123} +x_{124}\) \\[3pt]
\(s_{106}=\) & \(\displaystyle -1 +x_{120}^{2} +x_{123}^{2}\) \\
 & \(\displaystyle -x_{124}^{2}\) \\[3pt]
\(s_{107}=\) & \(\displaystyle -17 -8x_{4} -x_{4}^{2}\) \\
 & \(\displaystyle +2x_{4}x_{120} +2x_{4}x_{121} -2x_{4}x_{171}\) \\
 & \(\displaystyle +8x_{120} +2x_{120}x_{171} +8x_{121}\) \\
 & \(\displaystyle +2x_{121}x_{171} -x_{125}x_{126} -8x_{171}\) \\
 & \(\displaystyle -x_{171}^{2}\) \\[3pt]
\(s_{108}=\) & \(\displaystyle -x_{20} -x_{20}x_{214} -x_{20}x_{215}\) \\
 & \(\displaystyle +x_{126}\) \\[3pt]
\(s_{109}=\) & \(\displaystyle -4x_{128} -4x_{128}x_{133} +x_{134}\) \\[3pt]
\(s_{110}=\) & \(\displaystyle -x_{120}x_{128} -x_{121}x_{128} +x_{135}\) \\[3pt]
\(s_{111}=\) & \(\displaystyle -1 +x_{127}^{2} +x_{128}^{2}\) \\
 & \(\displaystyle -x_{135}^{2}\) \\[3pt]
\(s_{112}=\) & \(\displaystyle -x_{120}x_{130} -x_{121}x_{130} +x_{136}\) \\[3pt]
\(s_{113}=\) & \(\displaystyle -1 +x_{129}^{2} +x_{130}^{2}\) \\
 & \(\displaystyle -x_{136}^{2}\) \\[3pt]
\(s_{114}=\) & \(\displaystyle -x_{132} -x_{132}x_{134} +x_{137}\) \\[3pt]
\(s_{115}=\) & \(\displaystyle -1 +x_{131}^{2} +x_{132}^{2}\) \\
 & \(\displaystyle -x_{137}^{2}\) \\[3pt]
\(s_{116}=\) & \(\displaystyle -x_{128}^{2} +x_{139}\) \\[3pt]
\(s_{117}=\) & \(\displaystyle x_{130} -x_{138}x_{139} -x_{139}\) \\[3pt]
\(s_{118}=\) & \(\displaystyle 1 -x_{120} -x_{121}\) \\
 & \(\displaystyle -x_{129}x_{140} +x_{134}\) \\[3pt]
\(s_{119}=\) & \(\displaystyle -x_{127} -x_{129}x_{141} +x_{131}\) \\[3pt]
\(s_{120}=\) & \(\displaystyle -25 -x_{11} -13x_{12}\) \\
 & \(\displaystyle -4x_{22} -2x_{25} -x_{29}\) \\
 & \(\displaystyle -x_{31} -x_{33} -x_{35}\) \\
 & \(\displaystyle -2x_{39} -4x_{41} -2x_{42}\) \\
 & \(\displaystyle -x_{45} -2x_{67} -4x_{128}x_{142}\) \\
 & \(\displaystyle +x_{132}\) \\[3pt]
\(s_{121}=\) & \(\displaystyle -x_{143}x_{144} +x_{145}\) \\[3pt]
\(s_{122}=\) & \(\displaystyle -1 -2x_{120}x_{145} -2x_{121}x_{145}\) \\
 & \(\displaystyle +x_{143}^{2} +x_{144}^{2}\) \\[3pt]
\(s_{123}=\) & \(\displaystyle -1 -x_{3}^{2} +2x_{3}x_{120}\) \\
 & \(\displaystyle +2x_{3}x_{121} -x_{146}\) \\[3pt]
\(s_{124}=\) & \(\displaystyle -4 -2x_{2} -x_{2}x_{4}\) \\
 & \(\displaystyle +x_{3}x_{143} -2x_{4} -x_{5}\) \\
 & \(\displaystyle -x_{144} -x_{146}x_{147}\) \\[3pt]
\(s_{125}=\) & \(\displaystyle -x_{148}x_{149} +x_{150}\) \\[3pt]
\(s_{126}=\) & \(\displaystyle -1 -2x_{120}x_{150} -2x_{121}x_{150}\) \\
 & \(\displaystyle +x_{148}^{2} +x_{149}^{2}\) \\[3pt]
\(s_{127}=\) & \(\displaystyle -5 -4x_{4} -x_{4}^{2}\) \\
 & \(\displaystyle +2x_{4}x_{120} +2x_{4}x_{121} +4x_{120}\) \\
 & \(\displaystyle +4x_{121} -x_{151}\) \\[3pt]
\(s_{128}=\) & \(\displaystyle x_{4}x_{148} -x_{20} +2x_{148}\) \\
 & \(\displaystyle -x_{149} -x_{151}x_{152}\) \\[3pt]
\(s_{129}=\) & \(\displaystyle -1 -x_{20}^{2} +2x_{20}x_{120}\) \\
 & \(\displaystyle +2x_{20}x_{121} -x_{153}\) \\[3pt]
\(s_{130}=\) & \(\displaystyle -3 -x_{4} +x_{20}x_{148}\) \\
 & \(\displaystyle -x_{149} -x_{153}x_{154} -x_{171}\) \\[3pt]
\(s_{131}=\) & \(\displaystyle -10 -6x_{4} -x_{4}^{2}\) \\
 & \(\displaystyle +2x_{4}x_{120} +2x_{4}x_{121} -2x_{4}x_{171}\) \\
 & \(\displaystyle +6x_{120} +2x_{120}x_{171} +6x_{121}\) \\
 & \(\displaystyle +2x_{121}x_{171} -x_{155} -6x_{171}\) \\
 & \(\displaystyle -x_{171}^{2}\) \\[3pt]
\(s_{132}=\) & \(\displaystyle -1 +x_{4}x_{143} +3x_{143}\) \\
 & \(\displaystyle +x_{143}x_{171} -x_{144} -x_{155}x_{156}\) \\
 & \(\displaystyle -x_{169} -x_{170}\) \\[3pt]
\(s_{133}=\) & \(\displaystyle -x_{157}x_{158} +x_{159}\) \\[3pt]
\(s_{134}=\) & \(\displaystyle -1 -2x_{120}x_{159} -2x_{121}x_{159}\) \\
 & \(\displaystyle +x_{157}^{2} +x_{158}^{2}\) \\[3pt]
\(s_{135}=\) & \(\displaystyle -1 +x_{4}x_{157} -x_{155}x_{160}\) \\
 & \(\displaystyle +3x_{157} +x_{157}x_{171} -x_{158}\) \\
 & \(\displaystyle -x_{201} -x_{202}\) \\[3pt]
\(s_{136}=\) & \(\displaystyle -x_{161}x_{162} +x_{163}\) \\[3pt]
\(s_{137}=\) & \(\displaystyle -1 -2x_{120}x_{163} -2x_{121}x_{163}\) \\
 & \(\displaystyle +x_{161}^{2} +x_{162}^{2}\) \\[3pt]
\(s_{138}=\) & \(\displaystyle -1 +x_{4}x_{161} -x_{155}x_{164}\) \\
 & \(\displaystyle +3x_{161} +x_{161}x_{171} -x_{162}\) \\
 & \(\displaystyle -x_{214} -x_{215}\) \\[3pt]
\(s_{139}=\) & \(\displaystyle -x_{165}x_{166} +x_{167}\) \\[3pt]
\(s_{140}=\) & \(\displaystyle -1 -2x_{120}x_{167} -2x_{121}x_{167}\) \\
 & \(\displaystyle +x_{165}^{2} +x_{166}^{2}\) \\[3pt]
\(s_{141}=\) & \(\displaystyle -2 -x_{4} +x_{4}x_{165}\) \\
 & \(\displaystyle -x_{4}x_{169} -x_{4}x_{170} +4x_{165}\) \\
 & \(\displaystyle +x_{165}x_{171} -x_{166} -x_{168}\) \\
 & \(\displaystyle -x_{168}x_{176} -x_{168}x_{177} -3x_{169}\) \\
 & \(\displaystyle -2x_{170}\) \\[3pt]
\(s_{142}=\) & \(\displaystyle -x_{172}x_{173} +x_{174}\) \\[3pt]
\(s_{143}=\) & \(\displaystyle -1 -2x_{120}x_{174} -2x_{121}x_{174}\) \\
 & \(\displaystyle +x_{172}^{2} +x_{173}^{2}\) \\[3pt]
\(s_{144}=\) & \(\displaystyle -x_{20} -x_{20}x_{176} -x_{20}x_{177}\) \\
 & \(\displaystyle +x_{178}\) \\[3pt]
\(s_{145}=\) & \(\displaystyle x_{4}x_{172} +x_{171}x_{172} +4x_{172}\) \\
 & \(\displaystyle -x_{173} -x_{175}x_{178} -x_{176}\) \\[3pt]
\(s_{146}=\) & \(\displaystyle -x_{179}x_{180} +x_{181}\) \\[3pt]
\(s_{147}=\) & \(\displaystyle -1 -2x_{120}x_{181} -2x_{121}x_{181}\) \\
 & \(\displaystyle +x_{179}^{2} +x_{180}^{2}\) \\[3pt]
\(s_{148}=\) & \(\displaystyle x_{4}x_{179} +x_{171}x_{179} -x_{178}x_{182}\) \\
 & \(\displaystyle +4x_{179} -x_{180} -x_{183}\) \\[3pt]
\(s_{149}=\) & \(\displaystyle x_{176} +x_{177} -x_{183}\) \\
 & \(\displaystyle -x_{184}\) \\[3pt]
\(s_{150}=\) & \(\displaystyle -x_{185}x_{186} +x_{187}\) \\[3pt]
\(s_{151}=\) & \(\displaystyle -1 -2x_{120}x_{187} -2x_{121}x_{187}\) \\
 & \(\displaystyle +x_{185}^{2} +x_{186}^{2}\) \\[3pt]
\(s_{152}=\) & \(\displaystyle x_{4}x_{185} +x_{171}x_{185} -x_{178}x_{188}\) \\
 & \(\displaystyle +4x_{185} -x_{186} -x_{189}\) \\[3pt]
\(s_{153}=\) & \(\displaystyle x_{176} +x_{177} -x_{189}\) \\
 & \(\displaystyle -x_{190}\) \\[3pt]
\(s_{154}=\) & \(\displaystyle -x_{191}x_{192} +x_{193}\) \\[3pt]
\(s_{155}=\) & \(\displaystyle -1 -2x_{120}x_{193} -2x_{121}x_{193}\) \\
 & \(\displaystyle +x_{191}^{2} +x_{192}^{2}\) \\[3pt]
\(s_{156}=\) & \(\displaystyle x_{4}x_{191} +x_{171}x_{191} -x_{178}x_{194}\) \\
 & \(\displaystyle +4x_{191} -x_{192} -x_{195}\) \\[3pt]
\(s_{157}=\) & \(\displaystyle x_{176} +x_{177} -x_{195}\) \\
 & \(\displaystyle -x_{196}\) \\[3pt]
\(s_{158}=\) & \(\displaystyle -x_{197}x_{198} +x_{199}\) \\[3pt]
\(s_{159}=\) & \(\displaystyle -1 -2x_{120}x_{199} -2x_{121}x_{199}\) \\
 & \(\displaystyle +x_{197}^{2} +x_{198}^{2}\) \\[3pt]
\(s_{160}=\) & \(\displaystyle -x_{20} -x_{20}x_{201} -x_{20}x_{202}\) \\
 & \(\displaystyle +x_{203}\) \\[3pt]
\(s_{161}=\) & \(\displaystyle x_{4}x_{197} +x_{171}x_{197} +4x_{197}\) \\
 & \(\displaystyle -x_{198} -x_{200}x_{203} -x_{201}\) \\[3pt]
\(s_{162}=\) & \(\displaystyle -x_{204}x_{205} +x_{206}\) \\[3pt]
\(s_{163}=\) & \(\displaystyle -1 -2x_{120}x_{206} -2x_{121}x_{206}\) \\
 & \(\displaystyle +x_{204}^{2} +x_{205}^{2}\) \\[3pt]
\(s_{164}=\) & \(\displaystyle x_{4}x_{204} +x_{171}x_{204} -x_{203}x_{207}\) \\
 & \(\displaystyle +4x_{204} -x_{205} -x_{208}\) \\[3pt]
\(s_{165}=\) & \(\displaystyle x_{201} +x_{202} -x_{208}\) \\
 & \(\displaystyle -x_{209}\) \\[3pt]
\(s_{166}=\) & \(\displaystyle -x_{210}x_{211} +x_{212}\) \\[3pt]
\(s_{167}=\) & \(\displaystyle -1 -2x_{120}x_{212} -2x_{121}x_{212}\) \\
 & \(\displaystyle +x_{210}^{2} +x_{211}^{2}\) \\[3pt]
\(s_{168}=\) & \(\displaystyle x_{4}x_{210} -x_{126}x_{213} +x_{171}x_{210}\) \\
 & \(\displaystyle +4x_{210} -x_{211} -x_{214}\) \\[3pt]
\(s_{169}=\) & \(\displaystyle -x_{216}x_{217} +x_{218}\) \\[3pt]
\(s_{170}=\) & \(\displaystyle -1 -2x_{120}x_{218} -2x_{121}x_{218}\) \\
 & \(\displaystyle +x_{216}^{2} +x_{217}^{2}\) \\[3pt]
\(s_{171}=\) & \(\displaystyle x_{4}x_{216} -x_{126}x_{219} +x_{171}x_{216}\) \\
 & \(\displaystyle +4x_{216} -x_{217} -x_{220}\) \\[3pt]
\(s_{172}=\) & \(\displaystyle x_{214} +x_{215} -x_{220}\) \\
 & \(\displaystyle -x_{221}\) \\[3pt]
\(s_{173}=\) & \(\displaystyle -x_{222}x_{223} +x_{224}\) \\[3pt]
\(s_{174}=\) & \(\displaystyle -1 -2x_{120}x_{224} -2x_{121}x_{224}\) \\
 & \(\displaystyle +x_{222}^{2} +x_{223}^{2}\) \\[3pt]
\(s_{175}=\) & \(\displaystyle -x_{20} -x_{20}x_{169} -x_{20}x_{170}\) \\
 & \(\displaystyle +x_{228}\) \\[3pt]
\(s_{176}=\) & \(\displaystyle x_{4}x_{222} +x_{171}x_{222} +4x_{222}\) \\
 & \(\displaystyle -x_{223} -x_{225}x_{228} -x_{226}\) \\[3pt]
\(s_{177}=\) & \(\displaystyle x_{169} +x_{170} -x_{226}\) \\
 & \(\displaystyle -x_{227}\) \\[3pt]
\(s_{178}=\) & \(\displaystyle x_{4}x_{128} +x_{127} +4x_{128}\) \\
 & \(\displaystyle +x_{128}x_{171} -x_{135} -x_{178}x_{229}\) \\
 & \(\displaystyle -x_{230}\) \\[3pt]
\(s_{179}=\) & \(\displaystyle x_{176} +x_{177} -x_{230}\) \\
 & \(\displaystyle -x_{231}\) \\[3pt]
\(s_{180}=\) & \(\displaystyle -1639 -176x_{0} -5x_{0}^{2}\) \\
 & \(\displaystyle -8x_{0}x_{1} -16x_{0}x_{2} -8x_{0}x_{8}\) \\
 & \(\displaystyle +2x_{0}x_{143} +4x_{0}x_{148} -154x_{1}\) \\
 & \(\displaystyle -5x_{1}^{2} -16x_{1}x_{2} -8x_{1}x_{8}\) \\
 & \(\displaystyle +4x_{1}x_{148} +2x_{1}x_{161} -312x_{2}\) \\
 & \(\displaystyle -17x_{2}^{2} -18x_{2}x_{8} +8x_{2}x_{148}\) \\
 & \(\displaystyle +2x_{2}x_{157} -x_{3}^{2} +2x_{3}x_{165}\) \\
 & \(\displaystyle -160x_{8} -5x_{8}^{2} +4x_{8}x_{148}\) \\
 & \(\displaystyle +2x_{8}x_{157} -3x_{11} -x_{11}^{2}\) \\
 & \(\displaystyle +2x_{11}x_{172} -15x_{12} -x_{12}^{2}\) \\
 & \(\displaystyle +2x_{12}x_{179} -2x_{13} -x_{13}^{2}\) \\
 & \(\displaystyle +2x_{13}x_{185} -2x_{14} -x_{14}^{2}\) \\
 & \(\displaystyle +2x_{14}x_{191} -2x_{15} -x_{15}^{2}\) \\
 & \(\displaystyle +2x_{15}x_{197} -2x_{16} -x_{16}^{2}\) \\
 & \(\displaystyle +2x_{16}x_{204} -2x_{17} -x_{17}^{2}\) \\
 & \(\displaystyle +2x_{17}x_{210} -2x_{18} -x_{18}^{2}\) \\
 & \(\displaystyle +2x_{18}x_{216} -2x_{19} -x_{19}^{2}\) \\
 & \(\displaystyle +2x_{19}x_{222} -4x_{22} -2x_{25}\) \\
 & \(\displaystyle -x_{29} -x_{31} -x_{33}\) \\
 & \(\displaystyle -x_{35} -2x_{39} -4x_{41}\) \\
 & \(\displaystyle -2x_{42} -x_{45} -2x_{67}\) \\
 & \(\displaystyle -x_{120}x_{232} -x_{121}x_{232} +x_{128}\) \\
 & \(\displaystyle +24x_{143} -x_{143}^{2} +76x_{148}\) \\
 & \(\displaystyle -x_{148}^{2} +8x_{157} -x_{157}^{2}\) \\
 & \(\displaystyle +2x_{161} -x_{161}^{2} -x_{165}^{2}\) \\
 & \(\displaystyle +2x_{172} -x_{172}^{2} +2x_{179}\) \\
 & \(\displaystyle -x_{179}^{2} +2x_{185} -x_{185}^{2}\) \\
 & \(\displaystyle +2x_{191} -x_{191}^{2} +2x_{197}\) \\
 & \(\displaystyle -x_{197}^{2} +2x_{204} -x_{204}^{2}\) \\
 & \(\displaystyle +2x_{210} -x_{210}^{2} +2x_{216}\) \\
 & \(\displaystyle -x_{216}^{2} +2x_{222} -x_{222}^{2}\) \\
 & \(\displaystyle +x_{232}\) \\[3pt]
\(s_{181}=\) & \(\displaystyle -117 -3x_{0} -3x_{1}\) \\
 & \(\displaystyle -7x_{2} -x_{2}x_{4} -x_{3}\) \\
 & \(\displaystyle -4x_{4} -x_{5} -3x_{8}\) \\
 & \(\displaystyle -2x_{11} -14x_{12} -x_{13}\) \\
 & \(\displaystyle -x_{14} -x_{15} -x_{16}\) \\
 & \(\displaystyle -x_{17} -x_{18} -x_{19}\) \\
 & \(\displaystyle -x_{20} -4x_{22} -2x_{25}\) \\
 & \(\displaystyle -x_{29} -x_{31} -x_{33}\) \\
 & \(\displaystyle -x_{35} -2x_{39} -4x_{41}\) \\
 & \(\displaystyle -2x_{42} -x_{45} -2x_{67}\) \\
 & \(\displaystyle +x_{119} -x_{169} -x_{170}\) \\
 & \(\displaystyle -x_{171} -x_{201} -x_{202}\) \\
 & \(\displaystyle -x_{214} -x_{215}\) \\[3pt]
\end{longtable}
\endgroup

\section{The four fixed compact coefficients}
\label{app:constants}
The recipe below is reproduced from the pinned indexed native source.
It supplies all four constants in Section~\ref{sec:compact-unscaled};
the symbol \(K\) in this recipe is its \(K_0\). All bounds and
indices here are fixed integers, not additional unknowns. The code in
\code{compact/rebuild_coefficients.py} independently evaluates the
recipe and checks its four literal payloads.
% Exact finite constant recipe from compact/proof_sources/indexed_constants.tex.
\begingroup
\providecommand{\fl}[1]{\left\lfloor#1\right\rfloor}
\providecommand{\rg}[2]{\{#1,\ldots,#2\}}
\textbf{The four fixed integers.}
Every index, set, sum, product and floor below defines a constant, not
another unknown. Put
\[
\rho=2^{67},\quad B=\rho^{8807},\quad\beta=B-1,\quad
G_m(x)=\sum_{j=0}^{m-1}x^j.
\]
Define
\[
\begin{split}
\mathcal E={}&\rg0{942}\cup\rg{983}{992}\cup\rg{1003}{1011}\\
&\cup\rg{1013}{1031}\cup\{1063,1064,1095,2195,3296\}.
\end{split}
\]
Then
\begin{align*}
M_C={}&B-1-\sum_{e\in\mathcal E\setminus\{1\}}\rho^e-2\rho^{1063},\\
M_F={}&B-1+4+(2^{51}-2)(\rho^{4393}+\rho^{5495})\\
&+\rho^{1013}G_{10}(\rho)+\rho^{1065}G_9(\rho)
  +\rho^{2195}+\rho^{3296}\\
&+(2^{65}-1)\rho^{8792}G_{13}(\rho)+(2^{19}-1)\rho^{8805},\\
K={}&\rho^{30}+\rho^{62}+\rho^{1100}+\rho^{2201}+\rho^{8849}\\
&+(\rho^{4393}+\rho^{5495})\Phi(1/\rho),
\end{align*}
where
\[
\begin{split}
\Phi(x)={}&(1+2^{51})G_{943}(x)
 -2^{10}\sum_{i=0}^{942}v(t_i)x^i\\
&+\sum_{(b,e,m)\in\mathcal C}2^b x^e G_m(2x),\\
\mathcal C={}&\{(10,1023,8),(18,1013,9),(27,1003,8),\\
&(35,983,9),(44,1031,1),(45,1022,1),(46,1011,1),\\
&(47,992,1),(48,1095,1),(49,2195,1),(50,3296,1)\}.
\end{split}
\]
The rational intermediate in K cancels exactly; all four results are positive integers.

\textbf{The 943 indexed fixed tiles.}
Use the actions $u_j,v_j,d_j$ decoded below. Let
\[
\begin{gathered}
h(q,b)=2(q+1)+b,\quad
\mathcal R=\{u_j:d_j=1,\ 0\le j<430\},\\
\mathcal L=\{u_j:d_j=0,\ 0\le j<430\},\\
r_q=6+|\mathcal R\cap\rg0{q-1}|,\quad
l_q=205+|\mathcal L\cap\rg0{q-1}|.
\end{gathered}
\]
An interval ending at $-1$ is empty. The table defines $t_i$ for
$0\le i\le942$; all index ranges are fixed:
\[
\setlength{\arraycolsep}{1.4pt}\begin{array}{c|l}
\text{index }i&\text{tile }t_i=(w,n,e,s)\\\hline
0,1&(2,435,3,1),\ (5,0,0,0)\\
j+2\ (j=2q+b<430)&((1-d_j)l_{u_j},h(q,b),d_jr_{u_j},v_j)\\
432+b\ (b=0,1)&(0,h(215,b),0,h(215,b))\\
434+b\ (b=0,1)&(0,b,0,b)\\
424+2r_q+b\ (q\in\mathcal R,b=0,1)&(r_q,b,0,h(q,b))\\
424+2l_q+b\ (q\in\mathcal L,b=0,1)&(0,b,l_q,h(q,b))\\
924+j\ (0\le j<4)&(0,j+430\fl{j/2},0,438)\\
928&(2,434,2,0)\\
929&(3,436,4,4)\\
930&(4,436,4,0)\\
497+t\ (t=434,435,436)&(1,438,1,t)\\
\substack{934+3a+b-6\fl{a/4}\\a\in\{0,1,4\},\ b=0,1,2}&(a,437,b,437)
\end{array}
\]
For $t=(w,n,e,s)$ put
\[
v(t)=w+2^8n+2^{17}e+2^{25}s+127\,2^{34}.
\]

\textbf{A closed decoder for the 430 actions.}
Put $D_b(A,k)=\fl{A/b^k}\bmod b$, where
$a\bmod b=a-b\fl{a/b}$. For $0\le q<215$, $b=0,1$, $j=2q+b$:
\[
\begin{gathered}
N=\mathcal A\bmod2^{636},\quad
T_0=\fl{\mathcal A/2^{636}}\bmod2^{1675},\\
E_0=\fl{\mathcal A/2^{2311}},\\
\nu_j=D_2(N,j),\quad \sigma_j=\sum_{k=0}^{j-1}\nu_k,\\
m_q=D_2(N,430+q),\quad \tau_q=\sum_{k=0}^{q-1}m_k,\\
F_q=14(1-m_q)+m_qD_{16}(E_0,\tau_q),\\
u_j=\nu_j(1+\sigma_j)+(1-\nu_j)D_{216}(T_0,j-\sigma_j),\\
v_j=D_2(F_q,2b),\quad d_j=D_2(F_q,2b+1).
\end{gathered}
\]
Empty sums are zero. Here $(v_0,d_0,u_0)=(1,1,1)$; the constructor tile
uses this actual first target, including its state renaming.

\textbf{The integer $\mathcal A$.}
The alphabet, in order of digit values $0,\ldots,61$, is\par
{\fontsize{9.2}{10.5}\selectfont\ttfamily 0123456789ABCDEFGHIJKLMNOPQRSTUVWXYZabcdefghijklmnopqrstuvwxyz}.\par
Concatenate the following 436 characters as digits $e_0,\ldots,e_{435}$;
put $\mathcal A=\sum_{j=0}^{435}e_j62^{435-j}$.\par
\begingroup\fontsize{9.2}{10.5}\selectfont\ttfamily
BxEnhvy6NipzkWmROo0GsA3uK3twWLxscLq20UTumy5Cdu461LwD6g7RDNzF8Y\par
iYDFTpB8ehAypd7d6OsOL4APamnJeNXQypUezCmW9n4oXP6XfgmiWuAWZ8ZJ7U\par
OP7gu4kVdlGZqZHULQMar1dXNeJW4i7oNVI8usp8BlovjjZbwPCovZI328fRUq\par
4q92R81s5GjWQaQfGhkq9UYMK3rRMaBmsOIBkJqc2cj2VRIZ6mjzkBIOkW9SjZ\par
FzeA6J5u4j3sBL3SdgwXzM9mHvQivu8EJ1TBTpRhHDqZIbmYgwQ8VqTAFotl6F\par
We1YC554GKrVHsA8oUGKSA6F6k5KREBSD0Tgd7bFALY84LLbpw574BQyvdoHrl\par
0Oe35uwDf15DpT6VJWZBjSbHPpPAINYTte0eUxnTFTKBlGj8Am6sb1gLldh8WZ\par
8P\par
\endgroup
\endgroup

\end{document}
